
\Data\ID{1003028801}
\Updated{2020-07-15 03:07:12}
\Level{A}
\SubArea{Higher degree equations and inequalities}
\LANGen{
\Question{
Find the sum of all the natural roots of the next equation. Multiple roots, if any, count only once.
\[ \left(x^2+4\right)\left(x^2-9\right)=0 \]}
{\( 3 \)}{\( -5 \)}{\( 5 \)}{\( -1 \)}{}{}}
\LANGpl{
\Question{
Znajdź sumę wszystkich naturalnych pierwiastków następującego równania. Wielokrotne pierwiastki, jeśli takie istnieją, liczy się tylko raz. 
\[ \left(x^2+4\right)\left(x^2-9\right)=0 \]}
{\( 3 \)}{\( -5 \)}{\( 5 \)}{\( -1 \)}{}{}}
\LANGcs{
\Question{
Určete součet všech přirozených kořenů dané rovnice. Případné vícenásobné kořeny započítejte jen jednou.
\[ \left(x^2+4\right)\left(x^2-9\right)=0 \]}
{\( 3 \)}{\( -5 \)}{\( 5 \)}{\( -1 \)}{}{}}
\LANGsk{
\Question{
Určte súčet všetkých prirodzených koreňov danej rovnice. Prípadné viacnásobné korene započítajte len raz.
\[ \left(x^2+4\right)\left(x^2-9\right)=0 \]}
{\( 3 \)}{\( -5 \)}{\( 5 \)}{\( -1 \)}{}{}}
\LANGes{
\Question{
Calcula la suma de todas las raíces naturales de la siguiente ecuación. Si hay raíces múltiples, cuéntalas solo una vez.
\[ \left(x^2+4\right)\left(x^2-9\right)=0 \]}
{\( 3 \)}{\( -5 \)}{\( 5 \)}{\( -1 \)}{}{}}
\End

\Data\ID{1003028802}
\Updated{2020-07-15 03:07:12}
\Level{A}
\SubArea{Higher degree equations and inequalities}
\LANGen{
\Question{
Find the sum of all the real roots of the next equation. Multiple roots, if any, count only once.
\[ \left(x^2+9\right)\left(x^2-1\right)\left(x+1\right)=0 \]}
{\( 0 \)}{\( 1 \)}{\( -1 \)}{\( 3 \)}{}{}}
\LANGpl{
\Question{
Znajdź sumę wszystkich rzeczywistych pierwiastków następnego równania. Wielokrotne pierwiastki, jeśli takie istnieją, liczy się tylko raz.
\[ \left(x^2+9\right)\left(x^2-1\right)\left(x+1\right)=0 \]}
{\( 0 \)}{\( 1 \)}{\( -1 \)}{\( 3 \)}{}{}}
\LANGcs{
\Question{
Určete součet všech reálných kořenů dané rovnice. Případné vícenásobné kořeny započítejte jen jednou.
\[ \left(x^2+9\right)\left(x^2-1\right)\left(x+1\right)=0 \]}
{\( 0 \)}{\( 1 \)}{\( -1 \)}{\( 3 \)}{}{}}
\LANGsk{
\Question{
Určte súčet všetkých reálnych koreňov danej rovnice. Prípadné viacnásobné korene započítajte len raz.
\[ \left(x^2+9\right)\left(x^2-1\right)\left(x+1\right)=0 \]}
{\( 0 \)}{\( 1 \)}{\( -1 \)}{\( 3 \)}{}{}}
\LANGes{
\Question{
Calcula la suma de todas las raíces reales de la siguiente ecuación. Si hay raíces múltiples, cuéntalas solo una vez.
\[ \left(x^2+9\right)\left(x^2-1\right)\left(x+1\right)=0 \]}
{\( 0 \)}{\( 1 \)}{\( -1 \)}{\( 3 \)}{}{}}
\End

\Data\ID{1003028803}
\Updated{2020-07-15 03:07:12}
\Level{A}
\SubArea{Higher degree equations and inequalities}
\LANGen{
\Question{
Find all the real roots of the equation.
\[ \left(x^2+1\right)\left(x^2-2\right)\left(x+2\right)=0 \]}
{\( \left\{-2;-\sqrt2;\sqrt2\right\} \)}{\( \left\{-2;-\sqrt2;-1;1;\sqrt2\right\} \)}{\( \left\{-2;-1;1;2\right\} \)}{\( \left\{-2;-\sqrt2;-1;\sqrt2\right\} \)}{}{}}
\LANGpl{
\Question{
Wyznacz wszystkie rzeczywiste pierwiastki równania.
\[ \left(x^2+1\right)\left(x^2-2\right)\left(x+2\right)=0 \]}
{\( \left\{-2;-\sqrt2;\sqrt2\right\} \)}{\( \left\{-2;-\sqrt2;-1;1;\sqrt2\right\} \)}{\( \left\{-2;-1;1;2\right\} \)}{\( \left\{-2;-\sqrt2;-1;\sqrt2\right\} \)}{}{}}
\LANGcs{
\Question{
Určete množinu všech reálných kořenů rovnice.
\[ \left(x^2+1\right)\left(x^2-2\right)\left(x+2\right)=0 \]}
{\( \left\{-2;-\sqrt2;\sqrt2\right\} \)}{\( \left\{-2;-\sqrt2;-1;1;\sqrt2\right\} \)}{\( \left\{-2;-1;1;2\right\} \)}{\( \left\{-2;-\sqrt2;-1;\sqrt2\right\} \)}{}{}}
\LANGsk{
\Question{
Určte množinu všetkých reálnych koreňov rovnice.
\[ \left(x^2+1\right)\left(x^2-2\right)\left(x+2\right)=0 \]}
{\( \left\{-2;-\sqrt2;\sqrt2\right\} \)}{\( \left\{-2;-\sqrt2;-1;1;\sqrt2\right\} \)}{\( \left\{-2;-1;1;2\right\} \)}{\( \left\{-2;-\sqrt2;-1;\sqrt2\right\} \)}{}{}}
\LANGes{
\Question{
Determina todas las raíces reales de la ecuación.
\[ \left(x^2+1\right)\left(x^2-2\right)\left(x+2\right)=0 \]}
{\( \left\{-2;-\sqrt2;\sqrt2\right\} \)}{\( \left\{-2;-\sqrt2;-1;1;\sqrt2\right\} \)}{\( \left\{-2;-1;1;2\right\} \)}{\( \left\{-2;-\sqrt2;-1;\sqrt2\right\} \)}{}{}}
\End

\Data\ID{1003028804}
\Updated{2020-07-15 03:07:12}
\Level{A}
\SubArea{Higher degree equations and inequalities}
\LANGen{
\Question{
Find the sum of all the natural roots of the next equation. Multiple roots, if any, count only once.
\[ \left(x^3+1\right)\left(x^2-4\right)=0 \]}
{\( 2 \)}{\( -1 \)}{\( 3 \)}{\( 1 \)}{}{}}
\LANGpl{
\Question{
Znajdź sumę wszystkich naturalnych pierwiastków następnego równania. Wielokrotne pierwiastki, jeśli takie istnieją, liczy się tylko raz.
\[ \left(x^3+1\right)\left(x^2-4\right)=0 \]}
{\( 2 \)}{\( -1 \)}{\( 3 \)}{\( 1 \)}{}{}}
\LANGcs{
\Question{
Určete součet všech přirozených kořenů dané rovnice. Případné vícenásobné kořeny započítejte jen jednou.
\[ \left(x^3+1\right)\left(x^2-4\right)=0 \]}
{\( 2 \)}{\( -1 \)}{\( 3 \)}{\( 1 \)}{}{}}
\LANGsk{
\Question{
Určte súčet všetkých prirodzených koreňov danej rovnice. Prípadné viacnásobné korene započítajte len raz.
\[ \left(x^3+1\right)\left(x^2-4\right)=0 \]}
{\( 2 \)}{\( -1 \)}{\( 3 \)}{\( 1 \)}{}{}}
\LANGes{
\Question{
Calcula la suma de todas las raíces naturales de la siguiente ecuación. Si hay raíces múltiples, cuéntalas solo una vez.
\[ \left(x^3+1\right)\left(x^2-4\right)=0 \]}
{\( 2 \)}{\( -1 \)}{\( 3 \)}{\( 1 \)}{}{}}
\End

\Data\ID{1003028805}
\Updated{2020-07-15 03:07:12}
\Level{A}
\SubArea{Higher degree equations and inequalities}
\LANGen{
\Question{
Find the sum of all the real roots of the next equation. Multiple roots, if any, count only once.
\[ \left(x^3+8\right)\left(x^2-1\right)\left(x+1\right)=0 \]}
{\( -2 \)}{\( 0 \)}{\( 2 \)}{\( -3 \)}{}{}}
\LANGpl{
\Question{
Znajdź sumę wszystkich rzeczywistych pierwiastków następnego równania. Wielokrotne pierwiastki, jeśli takie istnieją, liczy się tylko raz.
\[ \left(x^3+8\right)\left(x^2-1\right)\left(x+1\right)=0 \]}
{\( -2 \)}{\( 0 \)}{\( 2 \)}{\( -3 \)}{}{}}
\LANGcs{
\Question{
Určete součet všech reálných kořenů dané rovnice. Případné vícenásobné kořeny započítejte jen jednou.
\[ \left(x^3+8\right)\left(x^2-1\right)\left(x+1\right)=0 \]}
{\( -2 \)}{\( 0 \)}{\( 2 \)}{\( -3 \)}{}{}}
\LANGsk{
\Question{
Určte súčet všetkých reálnych koreňov danej rovnice. Prípadné viacnásobné korene započítajte len raz.
\[ \left(x^3+8\right)\left(x^2-1\right)\left(x+1\right)=0 \]}
{\( -2 \)}{\( 0 \)}{\( 2 \)}{\( -3 \)}{}{}}
\LANGes{
\Question{
Calcula la suma de todas las raíces reales de la siguiente ecuación. Si hay raíces múltiples, cuéntalas solo una vez.
\[ \left(x^3+8\right)\left(x^2-1\right)\left(x+1\right)=0 \]}
{\( -2 \)}{\( 0 \)}{\( 2 \)}{\( -3 \)}{}{}}
\End

\Data\ID{1003028806}
\Updated{2020-07-15 03:07:12}
\Level{A}
\SubArea{Higher degree equations and inequalities}
\LANGen{
\Question{
Find all the real roots of the equation.
\[ \left(x^4-1\right)\left(x^3-8\right)\left(x^2-2\right)=0 \]}
{\( \left\{-\sqrt2;-1;1;\sqrt2;2\right\} \)}{\( \left\{-2;-1;1;2\right\} \)}{\( \left\{-\sqrt2;-1;1;\sqrt2\right\} \)}{\( \left\{-\sqrt2;\sqrt2;2\right\} \)}{}{}}
\LANGpl{
\Question{
Wyznacz wszystkie rzeczywiste pierwiastki równania.
\[ \left(x^4-1\right)\left(x^3-8\right)\left(x^2-2\right)=0 \]}
{\( \left\{-\sqrt2;-1;1;\sqrt2;2\right\} \)}{\( \left\{-2;-1;1;2\right\} \)}{\( \left\{-\sqrt2;-1;1;\sqrt2\right\} \)}{\( \left\{-\sqrt2;\sqrt2;2\right\} \)}{}{}}
\LANGcs{
\Question{
Určete množinu všech reálných kořenů rovnice.
\[ \left(x^4-1\right)\left(x^3-8\right)\left(x^2-2\right)=0 \]}
{\( \left\{-\sqrt2;-1;1;\sqrt2;2\right\} \)}{\( \left\{-2;-1;1;2\right\} \)}{\( \left\{-\sqrt2;-1;1;\sqrt2\right\} \)}{\( \left\{-\sqrt2;\sqrt2;2\right\} \)}{}{}}
\LANGsk{
\Question{
Určte množinu všetkých reálnych koreňov rovnice.
\[ \left(x^4-1\right)\left(x^3-8\right)\left(x^2-2\right)=0 \]}
{\( \left\{-\sqrt2;-1;1;\sqrt2;2\right\} \)}{\( \left\{-2;-1;1;2\right\} \)}{\( \left\{-\sqrt2;-1;1;\sqrt2\right\} \)}{\( \left\{-\sqrt2;\sqrt2;2\right\} \)}{}{}}
\LANGes{
\Question{
Determina todas las raíces reales de la ecuación.
\[ \left(x^4-1\right)\left(x^3-8\right)\left(x^2-2\right)=0 \]}
{\( \left\{-\sqrt2;-1;1;\sqrt2;2\right\} \)}{\( \left\{-2;-1;1;2\right\} \)}{\( \left\{-\sqrt2;-1;1;\sqrt2\right\} \)}{\( \left\{-\sqrt2;\sqrt2;2\right\} \)}{}{}}
\End

\Data\ID{1003028901}
\Updated{2020-07-15 03:07:12}
\Level{A}
\SubArea{Higher degree equations and inequalities}
\LANGen{
\Question{
Find the factored form of the equation.
\[ x^4-16=0 \]}
{\( \left(x^2+4\right)\left(x-2\right)\left(x+2\right)=0 \)}{\( \left(x-2\right)^2\left(x+2\right)^2=0 \)}{\( \left(x-2\right)^4=0 \)}{\( \left(x+4\right)\left(x-2\right)\left(x+2\right)=0 \)}{}{}}
\LANGpl{
\Question{
Wyznacz wzór czynników równania.
\[ x^4-16=0 \]}
{\( \left(x^2+4\right)\left(x-2\right)\left(x+2\right)=0 \)}{\( \left(x-2\right)^2\left(x+2\right)^2=0 \)}{\( \left(x-2\right)^4=0 \)}{\( \left(x+4\right)\left(x-2\right)\left(x+2\right)=0 \)}{}{}}
\LANGcs{
\Question{
Určete součinový tvar rovnice.
\[ x^4-16=0 \]}
{\( \left(x^2+4\right)\left(x-2\right)\left(x+2\right)=0 \)}{\( \left(x-2\right)^2\left(x+2\right)^2=0 \)}{\( \left(x-2\right)^4=0 \)}{\( \left(x+4\right)\left(x-2\right)\left(x+2\right)=0 \)}{}{}}
\LANGsk{
\Question{
Určte súčinový tvar rovnice.
\[ x^4-16=0 \]}
{\( \left(x^2+4\right)\left(x-2\right)\left(x+2\right)=0 \)}{\( \left(x-2\right)^2\left(x+2\right)^2=0 \)}{\( \left(x-2\right)^4=0 \)}{\( \left(x+4\right)\left(x-2\right)\left(x+2\right)=0 \)}{}{}}
\LANGes{
\Question{
Determina la forma factorizada de la ecuación.
\[ x^4-16=0 \]}
{\( \left(x^2+4\right)\left(x-2\right)\left(x+2\right)=0 \)}{\( \left(x-2\right)^2\left(x+2\right)^2=0 \)}{\( \left(x-2\right)^4=0 \)}{\( \left(x+4\right)\left(x-2\right)\left(x+2\right)=0 \)}{}{}}
\End

\Data\ID{1003028902}
\Updated{2020-07-15 03:07:12}
\Level{A}
\SubArea{Higher degree equations and inequalities}
\LANGen{
\Question{
Find the factored form of the equation.
\[ x^4+2x^2+x^3+2x=0 \]}
{\( x\left(x^2+2\right)\left(x+1\right)=0 \)}{\( x\left(x^2+1\right)\left(x+1\right)=0 \)}{\( x\left(x^2+1\right)\left(x+2\right)=0 \)}{\( x\left(x^2+2\right)\left(x+2\right)=0 \)}{}{}}
\LANGpl{
\Question{
Wyznacz wzór czynników równania.
\[ x^4+2x^2+x^3+2x=0 \]}
{\( x\left(x^2+2\right)\left(x+1\right)=0 \)}{\( x\left(x^2+1\right)\left(x+1\right)=0 \)}{\( x\left(x^2+1\right)\left(x+2\right)=0 \)}{\( x\left(x^2+2\right)\left(x+2\right)=0 \)}{}{}}
\LANGcs{
\Question{
Určete součinový tvar rovnice.
\[ x^4+2x^2+x^3+2x=0 \]}
{\( x\left(x^2+2\right)\left(x+1\right)=0 \)}{\( x\left(x^2+1\right)\left(x+1\right)=0 \)}{\( x\left(x^2+1\right)\left(x+2\right)=0 \)}{\( x\left(x^2+2\right)\left(x+2\right)=0 \)}{}{}}
\LANGsk{
\Question{
Určte súčinový tvar rovnice.
\[ x^4+2x^2+x^3+2x=0 \]}
{\( x\left(x^2+2\right)\left(x+1\right)=0 \)}{\( x\left(x^2+1\right)\left(x+1\right)=0 \)}{\( x\left(x^2+1\right)\left(x+2\right)=0 \)}{\( x\left(x^2+2\right)\left(x+2\right)=0 \)}{}{}}
\LANGes{
\Question{
Determina la forma factorizada de la ecuación.
\[ x^4+2x^2+x^3+2x=0 \]}
{\( x\left(x^2+2\right)\left(x+1\right)=0 \)}{\( x\left(x^2+1\right)\left(x+1\right)=0 \)}{\( x\left(x^2+1\right)\left(x+2\right)=0 \)}{\( x\left(x^2+2\right)\left(x+2\right)=0 \)}{}{}}
\End

\Data\ID{1003028903}
\Updated{2020-07-15 03:07:12}
\Level{A}
\SubArea{Higher degree equations and inequalities}
\LANGen{
\Question{
Find the factored form of the equation.
\[ x^3-2x^2-x+2=0 \]}
{\( (x-2)(x-1)(x+1)=0 \)}{\( x(x-1)(x-2)=0 \)}{\( x^2(x-2)=0 \)}{\( x(x+1)(x-2)=0 \)}{}{}}
\LANGpl{
\Question{
Wyznacz wzór czynników równania.
\[ x^3-2x^2-x+2=0 \]}
{\( (x-2)(x-1)(x+1)=0 \)}{\( x(x-1)(x-2)=0 \)}{\( x^2(x-2)=0 \)}{\( x(x+1)(x-2)=0 \)}{}{}}
\LANGcs{
\Question{
Určete součinový tvar rovnice.
\[ x^3-2x^2-x+2=0 \]}
{\( (x-2)(x-1)(x+1)=0 \)}{\( x(x-1)(x-2)=0 \)}{\( x^2(x-2)=0 \)}{\( x(x+1)(x-2)=0 \)}{}{}}
\LANGsk{
\Question{
Určte súčinový tvar rovnice.
\[ x^3-2x^2-x+2=0 \]}
{\( (x-2)(x-1)(x+1)=0 \)}{\( x(x-1)(x-2)=0 \)}{\( x^2(x-2)=0 \)}{\( x(x+1)(x-2)=0 \)}{}{}}
\LANGes{
\Question{
Determina la forma factorizada de la ecuación.
\[ x^3-2x^2-x+2=0 \]}
{\( (x-2)(x-1)(x+1)=0 \)}{\( x(x-1)(x-2)=0 \)}{\( x^2(x-2)=0 \)}{\( x(x+1)(x-2)=0 \)}{}{}}
\End

\Data\ID{1003028904}
\Updated{2020-07-15 03:07:12}
\Level{A}
\SubArea{Higher degree equations and inequalities}
\LANGen{
\Question{
Find the factored form of the equation.
\[ x^4+2x^2=0 \]}
{\( x^2\left(x^2+2\right)=0 \)}{\( \left(x^2-2\right)\left(x^2+2\right)=0 \)}{\( \left(x^2+\sqrt2x\right)\left(x^2-\sqrt2x\right)=0 \)}{\( \left(x^2+\sqrt2x\right)\left(x^2+\sqrt2x\right)=0 \)}{}{}}
\LANGpl{
\Question{
Wyznacz wzór czynników równania.
\[ x^4+2x^2=0 \]}
{\( x^2\left(x^2+2\right)=0 \)}{\( \left(x^2-2\right)\left(x^2+2\right)=0 \)}{\( \left(x^2+\sqrt2x\right)\left(x^2-\sqrt2x\right)=0 \)}{\( \left(x^2+\sqrt2x\right)\left(x^2+\sqrt2x\right)=0 \)}{}{}}
\LANGcs{
\Question{
Určete součinový tvar rovnice.
\[ x^4+2x^2=0 \]}
{\( x^2\left(x^2+2\right)=0 \)}{\( \left(x^2-2\right)\left(x^2+2\right)=0 \)}{\( \left(x^2+\sqrt2x\right)\left(x^2-\sqrt2x\right)=0 \)}{\( \left(x^2+\sqrt2x\right)\left(x^2+\sqrt2x\right)=0 \)}{}{}}
\LANGsk{
\Question{
Určte súčinový tvar rovnice. 
\[ x^4+2x^2=0 \]}
{\( x^2\left(x^2+2\right)=0 \)}{\( \left(x^2-2\right)\left(x^2+2\right)=0 \)}{\( \left(x^2+\sqrt2x\right)\left(x^2-\sqrt2x\right)=0 \)}{\( \left(x^2+\sqrt2x\right)\left(x^2+\sqrt2x\right)=0 \)}{}{}}
\LANGes{
\Question{
Determina la forma factorizada de la ecuación.
\[ x^4+2x^2=0 \]}
{\( x^2\left(x^2+2\right)=0 \)}{\( \left(x^2-2\right)\left(x^2+2\right)=0 \)}{\( \left(x^2+\sqrt2x\right)\left(x^2-\sqrt2x\right)=0 \)}{\( \left(x^2+\sqrt2x\right)\left(x^2+\sqrt2x\right)=0 \)}{}{}}
\End

\Data\ID{1003028905}
\Updated{2020-07-15 03:07:12}
\Level{A}
\SubArea{Higher degree equations and inequalities}
\LANGen{
\Question{
Find the factored form of the equation.
\[ x^4+2x^3-3x-6=0 \]}
{\( \left(x^3-3\right)(x+2) = 0 \)}{\( \left(x^3+3\right)(x-2) = 0 \)}{\( \left(x^3-2\right)(x+3) = 0 \)}{\( \left(x^3+2\right)(x-3) = 0 \)}{}{}}
\LANGpl{
\Question{
Wyznacz wzór czynników równania.
\[ x^4+2x^3-3x-6=0 \]}
{\( \left(x^3-3\right)(x+2) = 0 \)}{\( \left(x^3+3\right)(x-2) = 0 \)}{\( \left(x^3-2\right)(x+3) = 0 \)}{\( \left(x^3+2\right)(x-3) = 0 \)}{}{}}
\LANGcs{
\Question{
Určete součinový tvar rovnice.
\[ x^4+2x^3-3x-6=0 \]}
{\( \left(x^3-3\right)(x+2) = 0 \)}{\( \left(x^3+3\right)(x-2) = 0 \)}{\( \left(x^3-2\right)(x+3) = 0 \)}{\( \left(x^3+2\right)(x-3) = 0 \)}{}{}}
\LANGsk{
\Question{
Určte súčinový tvar rovnice.
\[ x^4+2x^3-3x-6=0 \]}
{\( \left(x^3-3\right)(x+2) = 0 \)}{\( \left(x^3+3\right)(x-2) = 0 \)}{\( \left(x^3-2\right)(x+3) = 0 \)}{\( \left(x^3+2\right)(x-3) = 0 \)}{}{}}
\LANGes{
\Question{
Determina la forma factorizada de la ecuación.
\[ x^4+2x^3-3x-6=0 \]}
{\( \left(x^3-3\right)(x+2) = 0 \)}{\( \left(x^3+3\right)(x-2) = 0 \)}{\( \left(x^3-2\right)(x+3) = 0 \)}{\( \left(x^3+2\right)(x-3) = 0 \)}{}{}}
\End

\Data\ID{1003189104}
\Updated{2020-07-15 03:07:21}
\Level{A}
\SubArea{Higher degree equations and inequalities}
\LANGen{
\Question{
One of the solutions of the following equation equals \( 1 \). Find the sum of all the remaining solutions.
\[ x^4- x^3+ x - 1=0 \]}
{\( -1 \)}{\( 1 \)}{\( 0 \)}{\( 2 \)}{}{}}
\LANGpl{
\Question{
Jedno z rozwiązań podanego równania wynosi \( 1 \). Wyznacz sumę pozostałych rozwiązań. 
\[ x^4- x^3+ x - 1=0 \]}
{\( -1 \)}{\( 1 \)}{\( 0 \)}{\( 2 \)}{}{}}
\LANGcs{
\Question{
Jedno z řešení následující rovnice je $1$. Určete součet ostatních reálných řešení rovnice.
\[ x^4- x^3+ x - 1=0 \]}
{\( -1 \)}{\( 1 \)}{\( 0 \)}{\( 2 \)}{}{}}
\LANGsk{
\Question{
Jedno z riešení nasledujúcej rovnice je \( 1 \). Určte súčet ostatných reálnych riešení.
\[ x^4- x^3+ x - 1=0 \]}
{\( -1 \)}{\( 1 \)}{\( 0 \)}{\( 2 \)}{}{}}
\LANGes{
\Question{
Una de las soluciones de la siguiente ecuación es \( 1 \). Calcula la suma de todas las soluciones restantes.
\[ x^4- x^3+ x - 1=0 \]}
{\( -1 \)}{\( 1 \)}{\( 0 \)}{\( 2 \)}{}{}}
\End

\Data\ID{1003189105}
\Updated{2020-07-15 03:07:21}
\Level{A}
\SubArea{Higher degree equations and inequalities}
\LANGen{
\Question{
One of the roots of the following equation equals \( 2 \). Find the product of all the remaining roots.
\[ x^5- 2 x^4- 9 x + 18=0 \]}
{\( -3 \)}{\( -1 \)}{\( 1 \)}{\( 3 \)}{}{}}
\LANGpl{
\Question{
Jeden z pierwiastków podanego równania wynosi \( 2 \). Wyznacz iloczyn wszystkich pozostałych pierwiastków.
\[ x^5- 2 x^4- 9 x + 18=0 \]}
{\( -3 \)}{\( -1 \)}{\( 1 \)}{\( 3 \)}{}{}}
\LANGcs{
\Question{
Jeden z kořenů následující rovnice je $2$. Určete součin zbývajících reálných kořenů rovnice.
\[ x^5- 2 x^4- 9 x + 18=0 \]}
{\( -3 \)}{\( -1 \)}{\( 1 \)}{\( 3 \)}{}{}}
\LANGsk{
\Question{
Jeden z koreňov nasledujúcej rovnice je \( 2 \). Určte súčin zvyšných koreňov.
\[ x^5- 2 x^4- 9 x + 18=0 \]}
{\( -3 \)}{\( -1 \)}{\( 1 \)}{\( 3 \)}{}{}}
\LANGes{
\Question{
Una de las soluciones de la siguiente ecuación es \( 2 \). Calcula el producto de todas las raíces restantes.
\[ x^5- 2 x^4- 9 x + 18=0 \]}
{\( -3 \)}{\( -1 \)}{\( 1 \)}{\( 3 \)}{}{}}
\End

\Data\ID{1003189106}
\Updated{2020-07-15 03:07:21}
\Level{A}
\SubArea{Higher degree equations and inequalities}
\LANGen{
\Question{
Given \( x\in\mathbb{R} \), find the solution set of the following equation.
\[ 2x^4+x^3+2x^2+x=0 \]
Hint: Decompose the polynomial on the left side of the equation into factors.}
{\( \left\{-\frac12;0\right\} \)}{\( \{-1;0\} \)}{\( \left\{-1;0;\frac12;1\right\} \)}{\( \left\{-\frac{\sqrt2}2;0;1;\frac{\sqrt2}2\right\} \)}{}{}}
\LANGpl{
\Question{
Dany jest \( x\in\mathbb{R} \), wyznacz zbiór rozwiązań podanego równania. 
\[ 2x^4+x^3+2x^2+x=0 \]
Wskazówka: Rozłóż wielomian po lewej stronie równania na czynniki.}
{\( \left\{-\frac12;0\right\} \)}{\( \{-1;0\} \)}{\( \left\{-1;0;\frac12;1\right\} \)}{\( \left\{-\frac{\sqrt2}2;0;1;\frac{\sqrt2}2\right\} \)}{}{}}
\LANGcs{
\Question{
Najděte množinu řešení následující rovnice pro \( x\in\mathbb{R} \).
\[ 2x^4+x^3+2x^2+x=0 \]
Nápověda: Rozložte polynom na levé straně rovnice na součin kořenových činitelů.}
{\( \left\{-\frac12;0\right\} \)}{\( \{-1;0\} \)}{\( \left\{-1;0;\frac12;1\right\} \)}{\( \left\{-\frac{\sqrt2}2;0;1;\frac{\sqrt2}2\right\} \)}{}{}}
\LANGsk{
\Question{
Pre \( x\in\mathbb{R} \) určte množinu riešení nasledujúcej rovnice.
\[ 2x^4+x^3+2x^2+x=0 \]
Návod: Rozložte polynóm na ľavej strane rovnice na súčin koreňových činiteľov.}
{\( \left\{-\frac12;0\right\} \)}{\( \{-1;0\} \)}{\( \left\{-1;0;\frac12;1\right\} \)}{\( \left\{-\frac{\sqrt2}2;0;1;\frac{\sqrt2}2\right\} \)}{}{}}
\LANGes{
\Question{
Suponiendo que \( x\in\mathbb{R} \), calcula el conjunto de soluciones de la ecuación.
\[ 2x^4+x^3+2x^2+x=0 \]
Sugerencia: Factoriza el polinomio del miembro izquierdo de la ecuación en factores.}
{\( \left\{-\frac12;0\right\} \)}{\( \{-1;0\} \)}{\( \left\{-1;0;\frac12;1\right\} \)}{\( \left\{-\frac{\sqrt2}2;0;1;\frac{\sqrt2}2\right\} \)}{}{}}
\End

\Data\ID{1003189107}
\Updated{2020-07-15 03:07:21}
\Level{A}
\SubArea{Higher degree equations and inequalities}
\LANGen{
\Question{
Given \( x\in\mathbb{R} \), find the solution set of the following equation.
\[ -x^4+x^3-x^2+x=0 \]
Hint: Decompose the polynomial on the left side of the equation into factors.}
{\( \{ 0;1 \} \)}{\( \{-1;0\} \)}{\( \{-1;0;1\} \)}{\( \{-1;1\} \)}{}{}}
\LANGpl{
\Question{
Dany jest \( x\in\mathbb{R} \), wyznacz zbiór rozwiązań podanego równania.  
\[ -x^4+x^3-x^2+x=0 \]
Wskazówka: Rozłóż wielomian po lewej stronie równania na czynniki.}
{\( \{ 0;1 \} \)}{\( \{-1;0\} \)}{\( \{-1;0;1\} \)}{\( \{-1;1\} \)}{}{}}
\LANGcs{
\Question{
Najděte množinu řešení následující rovnice pro \( x\in\mathbb{R} \).
\[ -x^4+x^3-x^2+x=0 \]
Nápověda: Rozložte polynom na levé straně rovnice na součin kořenových činitelů.}
{\( \{ 0;1 \} \)}{\( \{-1;0\} \)}{\( \{-1;0;1\} \)}{\( \{-1;1\} \)}{}{}}
\LANGsk{
\Question{
Pre \( x\in\mathbb{R} \) určte množinu riešení nasledujúcej rovnice.
\[ -x^4+x^3-x^2+x=0 \]
Návod: Rozložte polynóm na ľavej strane rovnice na súčin koreňových činiteľov.}
{\( \{ 0;1 \} \)}{\( \{-1;0\} \)}{\( \{-1;0;1\} \)}{\( \{-1;1\} \)}{}{}}
\LANGes{
\Question{
Suponiendo que \( x\in\mathbb{R} \), calcula el conjunto de soluciones de la ecuación.
\[ -x^4+x^3-x^2+x=0 \]
Sugerencia:  Factoriza el polinomio del miembro izquierdo de la ecuación en factores.}
{\( \{ 0;1 \} \)}{\( \{-1;0\} \)}{\( \{-1;0;1\} \)}{\( \{-1;1\} \)}{}{}}
\End

\Data\ID{1003189108}
\Updated{2020-07-15 03:07:21}
\Level{A}
\SubArea{Higher degree equations and inequalities}
\LANGen{
\Question{
Find the product of all the real roots of the following equation.
\[ 5x^5-3x^3-5x^2+3=0 \]}
{\( -\frac35 \)}{\( \frac53 \)}{\( 0 \)}{\( 1 \)}{}{}}
\LANGpl{
\Question{
Wyznacz iloczyn wszystkich rzeczywistych pierwiastków podanego równania. 
\[ 5x^5-3x^3-5x^2+3=0 \]}
{\( -\frac35 \)}{\( \frac53 \)}{\( 0 \)}{\( 1 \)}{}{}}
\LANGcs{
\Question{
Určete součin všech reálných kořenů následující rovnice.
\[ 5x^5-3x^3-5x^2+3=0 \]}
{\( -\frac35 \)}{\( \frac53 \)}{\( 0 \)}{\( 1 \)}{}{}}
\LANGsk{
\Question{
Určte súčin všetkých reálnych koreňov nasledujúcej rovnice.
\[ 5x^5-3x^3-5x^2+3=0 \]}
{\( -\frac35 \)}{\( \frac53 \)}{\( 0 \)}{\( 1 \)}{}{}}
\LANGes{
\Question{
Calcula el producto de todas las raíces reales de la siguiente ecuación.
\[ 5x^5-3x^3-5x^2+3=0 \]}
{\( -\frac35 \)}{\( \frac53 \)}{\( 0 \)}{\( 1 \)}{}{}}
\End

\Data\ID{1003189109}
\Updated{2020-07-15 03:07:21}
\Level{A}
\SubArea{Higher degree equations and inequalities}
\LANGen{
\Question{
Find the sum of all the positive roots of the following equation.
\[ 16x^5-9x^3-16x^2+9=0 \]}
{\( \frac74 \)}{\( 1 \)}{\( \frac34 \)}{\( 0 \)}{}{}}
\LANGpl{
\Question{
Wyznacz sumę wszystkich dodatnich pierwiastków podanego równania. 
\[ 16x^5-9x^3-16x^2+9=0 \]}
{\( \frac74 \)}{\( 1 \)}{\( \frac34 \)}{\( 0 \)}{}{}}
\LANGcs{
\Question{
Určete součet všech kladných reálných kořenů následující rovnice.
\[ 16x^5-9x^3-16x^2+9=0 \]}
{\( \frac74 \)}{\( 1 \)}{\( \frac34 \)}{\( 0 \)}{}{}}
\LANGsk{
\Question{
Určte súčet všetkých kladných reálnych koreňov nasledujúcej rovnice.
\[ 16x^5-9x^3-16x^2+9=0 \]}
{\( \frac74 \)}{\( 1 \)}{\( \frac34 \)}{\( 0 \)}{}{}}
\LANGes{
\Question{
Calcula la suma de todas las raíces positivas de la siguiente ecuación.
\[ 16x^5-9x^3-16x^2+9=0 \]}
{\( \frac74 \)}{\( 1 \)}{\( \frac34 \)}{\( 0 \)}{}{}}
\End

\Data\ID{9000019801}
\Updated{2020-07-15 03:07:34}
\Level{A}
\SubArea{Higher degree equations and inequalities}
\LANGen{
\Question{
Assuming \(x\in \mathbb{N}\),
find the solution set of the following equation. 
\[ 
 x^{3} - 6x^{2} + 9x = 0
\]}
{\(\left \{3\right \}\)}{\(\emptyset \)}{\(\left \{0;3\right \}\)}{\(\left \{-3;3\right \}\)}{}{}}
\LANGpl{
\Question{
Zakładając, że \(x\in \mathbb{N}\),
wyznacz zbiór rozwiązań następującego równania.
\[ 
 x^{3} - 6x^{2} + 9x = 0
\]}
{\(\left \{3\right \}\)}{\(\emptyset \)}{\(\left \{0;3\right \}\)}{\(\left \{-3;3\right \}\)}{}{}}
\LANGcs{
\Question{
Určete množinu všech řešení dané rovnice v oboru přirozených čísel.
\[x^{3} - 6x^{2} + 9x = 0\]}
{\(\left \{3\right \}\)}{\(\emptyset \)}{\(\left \{0;3\right \}\)}{\(\left \{-3;3\right \}\)}{}{}}
\LANGsk{
\Question{
Určte množinu všetkých riešení danej rovnice pre \(x\in \mathbb{N}\). 
\[ 
 x^{3} - 6x^{2} + 9x = 0
\]}
{\(\left \{3\right \}\)}{\(\emptyset \)}{\(\left \{0;3\right \}\)}{\(\left \{-3;3\right \}\)}{}{}}
\LANGes{
\Question{
Suponiendo que \(x\in \mathbb{N}\), calcula el conjunto de soluciones de la siguiente ecuación.
\[ 
 x^{3} - 6x^{2} + 9x = 0
\]}
{\(\left \{3\right \}\)}{\(\emptyset \)}{\(\left \{0;3\right \}\)}{\(\left \{-3;3\right \}\)}{}{}}
\End

\Data\ID{9000019802}
\Updated{2020-07-15 03:07:34}
\Level{A}
\SubArea{Higher degree equations and inequalities}
\LANGen{
\Question{
Assuming \(x\in \mathbb{N}\),
find the solution set of the following equation. 
\[ 
 2x^{3} - 3x^{2} = 0
\]}
{\(\emptyset \)}{\(\left \{0\right \}\)}{\(\left \{2\right \}\)}{\(\left \{0; \frac{3} 
{2}\right \}\)}{}{}}
\LANGpl{
\Question{
Zakładając, że \(x\in \mathbb{N}\),
wyznacz zbiór rozwiązań następującego równania.
\[ 
 2x^{3} - 3x^{2} = 0
\]}
{\(\emptyset \)}{\(\left \{0\right \}\)}{\(\left \{2\right \}\)}{\(\left \{0; \frac{3} 
{2}\right \}\)}{}{}}
\LANGcs{
\Question{
Určete množinu všech řešení dané rovnice v oboru přirozených čísel.
\[2x^{3} - 3x^{2} = 0\]}
{\(\emptyset \)}{\(\left \{0\right \}\)}{\(\left \{2\right \}\)}{\(\left \{0; \frac{3} 
{2}\right \}\)}{}{}}
\LANGsk{
\Question{
Určte množinu všetkých riešení danej rovnice pre \(x\in \mathbb{N}\). 
\[ 
 2x^{3} - 3x^{2} = 0
\]}
{\(\emptyset \)}{\(\left \{0\right \}\)}{\(\left \{2\right \}\)}{\(\left \{0; \frac{3} 
{2}\right \}\)}{}{}}
\LANGes{
\Question{
Suponiendo que \(x\in \mathbb{N}\),
calcula el conjunto de soluciones de la siguiente ecuación. 
\[ 
 2x^{3} - 3x^{2} = 0
\]}
{\(\emptyset \)}{\(\left \{0\right \}\)}{\(\left \{2\right \}\)}{\(\left \{0; \frac{3} 
{2}\right \}\)}{}{}}
\End

\Data\ID{9000019807}
\Updated{2020-07-15 03:07:34}
\Level{A}
\SubArea{Higher degree equations and inequalities}
\LANGen{
\Question{
Assuming \(x\in \mathbb{R}\), find the solution set of the following equation. 
\[ 
 \left (3x + 2\right )\left (x\sqrt{2} + 1\right )\left (x^{2} + 1\right ) = 0
\]}
{\(\left \{-\frac{\sqrt{2}}
 {2} ;-\frac{2}
{3}\right \}\)}{\(\left \{-\frac{2}
{3}; \frac{1} 
{\sqrt{2}}\right \}\)}{\(\left \{\frac{2}
{3}; \frac{1} 
{\sqrt{2}}\right \}\)}{\(\left \{-1;-\frac{\sqrt{2}}
 {2} ;-\frac{2}
{3}\right \}\)}{}{}}
\LANGpl{
\Question{
Zakładając, że \(x\in \mathbb{R}\), wyznacz zbiór rozwiązań następującego równania.
\[ 
 \left (3x + 2\right )\left (x\sqrt{2} + 1\right )\left (x^{2} + 1\right ) = 0
\]}
{\(\left \{-\frac{\sqrt{2}}
 {2} ;-\frac{2}
{3}\right \}\)}{\(\left \{-\frac{2}
{3}; \frac{1} 
{\sqrt{2}}\right \}\)}{\(\left \{\frac{2}
{3}; \frac{1} 
{\sqrt{2}}\right \}\)}{\(\left \{-1;-\frac{\sqrt{2}}
 {2} ;-\frac{2}
{3}\right \}\)}{}{}}
\LANGcs{
\Question{
Určete množinu všech řešení dané rovnice.
\[\left (3x + 2\right )\left (x\sqrt{2} + 1\right )\left (x^{2} + 1\right ) = 0\]}
{\(\left \{-\frac{\sqrt{2}}
 {2} ;-\frac{2}
{3}\right \}\)}{\(\left \{-\frac{2}
{3}; \frac{1} 
{\sqrt{2}}\right \}\)}{\(\left \{\frac{2}
{3}; \frac{1} 
{\sqrt{2}}\right \}\)}{\(\left \{-1;-\frac{\sqrt{2}}
 {2} ;-\frac{2}
{3}\right \}\)}{}{}}
\LANGsk{
\Question{
Určte množinu všetkých riešení danej rovnice pre \(x\in \mathbb{R}\). 
\[ 
 \left (3x + 2\right )\left (x\sqrt{2} + 1\right )\left (x^{2} + 1\right ) = 0
\]}
{\(\left \{-\frac{\sqrt{2}}
 {2} ;-\frac{2}
{3}\right \}\)}{\(\left \{-\frac{2}
{3}; \frac{1} 
{\sqrt{2}}\right \}\)}{\(\left \{\frac{2}
{3}; \frac{1} 
{\sqrt{2}}\right \}\)}{\(\left \{-1;-\frac{\sqrt{2}}
 {2} ;-\frac{2}
{3}\right \}\)}{}{}}
\LANGes{
\Question{
Suponiendo que \(x\in \mathbb{R}\), calcula el conjunto de soluciones de la siguiente ecuación. 
\[ 
 \left (3x + 2\right )\left (x\sqrt{2} + 1\right )\left (x^{2} + 1\right ) = 0
\]}
{\(\left \{-\frac{\sqrt{2}}
 {2} ;-\frac{2}
{3}\right \}\)}{\(\left \{-\frac{2}
{3}; \frac{1} 
{\sqrt{2}}\right \}\)}{\(\left \{\frac{2}
{3}; \frac{1} 
{\sqrt{2}}\right \}\)}{\(\left \{-1;-\frac{\sqrt{2}}
 {2} ;-\frac{2}
{3}\right \}\)}{}{}}
\End

\Data\ID{9000025805}
\Updated{2020-07-15 03:07:34}
\Level{A}
\SubArea{Higher degree equations and inequalities}
\LANGen{
\Question{
In the following list identify a true statement on the function
\(f\).
\[ 
 f(x) = (x + 1)(x + 2)(x - 3)
\]}
{\(f(x) < 0 \iff x\in (-\infty ;-2)\cup (-1;3)\)}{\(f(x) < 0 \iff x\in \left (-\infty ;-\frac{3}
{2}\right )\cup (1;3)\)}{\(f(x) < 0 \iff x\in \left (-\infty ;-\frac{3}
{2}\right )\cup (3;\infty )\)}{\(f(x) < 0 \iff x\in \left (-\frac{3}
{2};-1\right )\cup (3;\infty )\)}{}{}}
\LANGpl{
\Question{
Które z poniższych stwierdzeń dotyczących funkcji 
\(f\) jest prawdziwe?
\[ 
 f\colon y = (x + 1)(x + 2)(x - 3)
\]}
{\(f(x) < 0 \iff x\in (-\infty ;-2)\cup (-1;3)\)}{\(f(x) < 0 \iff x\in \left (-\infty ;-\frac{3}
{2}\right )\cup (1;3)\)}{\(f(x) < 0 \iff x\in \left (-\infty ;-\frac{3}
{2}\right )\cup (3;\infty )\)}{\(f(x) < 0 \iff x\in \left (-\frac{3}
{2};-1\right )\cup (3;\infty )\)}{}{}}
\LANGcs{
\Question{
Který z následujících výroků o funkci \(f\) je pravdivý?
\[f\colon y = (x + 1)(x + 2)(x - 3)\]}
{\(f(x) < 0 \iff x\in (-\infty ;-2)\cup (-1;3)\)}{\(f(x) < 0 \iff x\in \left (-\infty ;-\frac{3}
{2}\right )\cup (1;3)\)}{\(f(x) < 0 \iff x\in \left (-\infty ;-\frac{3}
{2}\right )\cup (3;\infty )\)}{\(f(x) < 0 \iff x\in \left (-\frac{3}
{2};-1\right )\cup (3;\infty )\)}{}{}}
\LANGsk{
\Question{
Ktorý z nasledujúcich výrokov o funkcii 
\(f\) je pravdivý?
\[ 
 f\colon y = (x + 1)(x + 2)(x - 3)
\]}
{\(f(x) < 0 \iff x\in (-\infty ;-2)\cup (-1;3)\)}{\(f(x) < 0 \iff x\in \left (-\infty ;-\frac{3}
{2}\right )\cup (1;3)\)}{\(f(x) < 0 \iff x\in \left (-\infty ;-\frac{3}
{2}\right )\cup (3;\infty )\)}{\(f(x) < 0 \iff x\in \left (-\frac{3}
{2};-1\right )\cup (3;\infty )\)}{}{}}
\LANGes{
\Question{
En la siguiente lista, identifica la declaración verdadera sobre la función
\(f\).
\[ 
 f(x) = (x + 1)(x + 2)(x - 3)
\]}
{\(f(x) < 0 \iff x\in (-\infty ;-2)\cup (-1;3)\)}{\(f(x) < 0 \iff x\in \left (-\infty ;-\frac{3}
{2}\right )\cup (1;3)\)}{\(f(x) < 0 \iff x\in \left (-\infty ;-\frac{3}
{2}\right )\cup (3;\infty )\)}{\(f(x) < 0 \iff x\in \left (-\frac{3}
{2};-1\right )\cup (3;\infty )\)}{}{}}
\End

\Data\ID{9000028306}
\Updated{2020-07-15 03:07:34}
\Level{A}
\SubArea{Higher degree equations and inequalities}
\LANGen{
\Question{
Find the sum of all real solutions of the following equation. 
\[ 
 \left (3 - x\right )\left (x^{2} - 4\right ) = 0
\]}
{\(3\)}{\(0\)}{\(2\)}{\(5\)}{}{}}
\LANGpl{
\Question{
Oblicz sumę wszystkich rzeczywistych rozwiązań podanego równania.
\[ 
 \left (3 - x\right )\left (x^{2} - 4\right ) = 0
\]}
{\(3\)}{\(0\)}{\(2\)}{\(5\)}{}{}}
\LANGcs{
\Question{
Určete součet všech reálných kořenů dané rovnice.
\[ 
 \left (3 - x\right )\left (x^{2} - 4\right ) = 0
\]}
{\(3\)}{\(0\)}{\(2\)}{\(5\)}{}{}}
\LANGsk{
\Question{
Určte súčet všetkých reálnych koreňov danej rovnice. 
\[ 
 \left (3 - x\right )\left (x^{2} - 4\right ) = 0
\]}
{\(3\)}{\(0\)}{\(2\)}{\(5\)}{}{}}
\LANGes{
\Question{
Calcula la suma de todas las soluciones reales de la siguiente ecuación. 
\[ 
 \left (3 - x\right )\left (x^{2} - 4\right ) = 0
\]}
{\(3\)}{\(0\)}{\(2\)}{\(5\)}{}{}}
\End

\Data\ID{1003029002}
\Updated{2020-07-15 03:07:12}
\Level{B}
\SubArea{Higher degree equations and inequalities}
\LANGen{
\Question{
Find the solution set of the inequality.
\[ \left(x^2-1\right)\left(x^2-3\right) > 0 \]}
{\( \left(-\infty;-\sqrt3\right)\cup(-1;1)\cup\left(\sqrt3;\infty\right) \)}{\( \left(-\sqrt3;-1\right)\cup\left(1;\sqrt3\right) \)}{\( \left(-\infty;-\sqrt3\right)\cup(1;\sqrt3)\cup\left(\sqrt3;\infty\right) \)}{\( \left(-\infty;1\right)\cup\left(\sqrt3;\infty\right) \)}{}{}}
\LANGpl{
\Question{
Wyznacz zbiór rozwiązań podanej nierówności.
\[ \left(x^2-1\right)\left(x^2-3\right) > 0 \]}
{\( \left(-\infty;-\sqrt3\right)\cup(-1;1)\cup\left(\sqrt3;\infty\right) \)}{\( \left(-\sqrt3;-1\right)\cup\left(1;\sqrt3\right) \)}{\( \left(-\infty;-\sqrt3\right)\cup(1;\sqrt3)\cup\left(\sqrt3;\infty\right) \)}{\( \left(-\infty;1\right)\cup\left(\sqrt3;\infty\right) \)}{}{}}
\LANGcs{
\Question{
Určete množinu řešení nerovnice.
\[ \left(x^2-1\right)\left(x^2-3\right) > 0 \]}
{\( \left(-\infty;-\sqrt3\right)\cup(-1;1)\cup\left(\sqrt3;\infty\right) \)}{\( \left(-\sqrt3;-1\right)\cup\left(1;\sqrt3\right) \)}{\( \left(-\infty;-\sqrt3\right)\cup(1;\sqrt3)\cup\left(\sqrt3;\infty\right) \)}{\( \left(-\infty;1\right)\cup\left(\sqrt3;\infty\right) \)}{}{}}
\LANGsk{
\Question{
Určte množinu riešení nerovnice.
\[ \left(x^2-1\right)\left(x^2-3\right) > 0 \]}
{\( \left(-\infty;-\sqrt3\right)\cup(-1;1)\cup\left(\sqrt3;\infty\right) \)}{\( \left(-\sqrt3;-1\right)\cup\left(1;\sqrt3\right) \)}{\( \left(-\infty;-\sqrt3\right)\cup(1;\sqrt3)\cup\left(\sqrt3;\infty\right) \)}{\( \left(-\infty;1\right)\cup\left(\sqrt3;\infty\right) \)}{}{}}
\LANGes{
\Question{
Calcula el conjunto de soluciones de la inecuación.
\[ \left(x^2-1\right)\left(x^2-3\right) > 0 \]}
{\( \left(-\infty;-\sqrt3\right)\cup(-1;1)\cup\left(\sqrt3;\infty\right) \)}{\( \left(-\sqrt3;-1\right)\cup\left(1;\sqrt3\right) \)}{\( \left(-\infty;-\sqrt3\right)\cup(1;\sqrt3)\cup\left(\sqrt3;\infty\right) \)}{\( \left(-\infty;1\right)\cup\left(\sqrt3;\infty\right) \)}{}{}}
\End

\Data\ID{1003029003}
\Updated{2020-07-15 03:07:12}
\Level{B}
\SubArea{Higher degree equations and inequalities}
\LANGen{
\Question{
Find the solution set of the inequality.
\[ \left(x^3+8\right)\left(x^4+1\right)\left(x-3\right) < 0 \]}
{\( (-2;3) \)}{\( (-\infty;-2)\cup(3;\infty) \)}{\( (-2;\infty) \)}{\( (-\infty;-2)\cup(-1;1)\cup(3;\infty) \)}{}{}}
\LANGpl{
\Question{
Wyznacz zbiór rozwiązań podanej nierówności.
\[ \left(x^3+8\right)\left(x^4+1\right)\left(x-3\right) < 0 \]}
{\( (-2;3) \)}{\( (-\infty;-2)\cup(3;\infty) \)}{\( (-2;\infty) \)}{\( (-\infty;-2)\cup(-1;1)\cup(3;\infty) \)}{}{}}
\LANGcs{
\Question{
Určete množinu řešení nerovnice.
\[ \left(x^3+8\right)\left(x^4+1\right)\left(x-3\right) < 0 \]}
{\( (-2;3) \)}{\( (-\infty;-2)\cup(3;\infty) \)}{\( (-2;\infty) \)}{\( (-\infty;-2)\cup(-1;1)\cup(3;\infty) \)}{}{}}
\LANGsk{
\Question{
Určte množinu riešení nerovnice.
\[ \left(x^3+8\right)\left(x^4+1\right)\left(x-3\right) < 0 \]}
{\( (-2;3) \)}{\( (-\infty;-2)\cup(3;\infty) \)}{\( (-2;\infty) \)}{\( (-\infty;-2)\cup(-1;1)\cup(3;\infty) \)}{}{}}
\LANGes{
\Question{
Calcula el conjunto de soluciones de la inecuación.
\[ \left(x^3+8\right)\left(x^4+1\right)\left(x-3\right) < 0 \]}
{\( (-2;3) \)}{\( (-\infty;-2)\cup(3;\infty) \)}{\( (-2;\infty) \)}{\( (-\infty;-2)\cup(-1;1)\cup(3;\infty) \)}{}{}}
\End

\Data\ID{1003029004}
\Updated{2020-07-15 03:07:12}
\Level{B}
\SubArea{Higher degree equations and inequalities}
\LANGen{
\Question{
Find the solution set of the inequality.
\[ x^4-1\geq0 \]}
{\( (-\infty;-1]\cup[1;\infty) \)}{\( [-1;1] \)}{\( \mathbb{R} \)}{\( [ 1;\infty ) \)}{}{}}
\LANGpl{
\Question{
Wyznacz zbiór rozwiązań podanej nierówności.
\[ x^4-1\geq0 \]}
{\( (-\infty;-1\rangle\cup\langle1;\infty) \)}{\( \langle-1;1\rangle \)}{\( \mathbb{R} \)}{\( \langle 1;\infty ) \)}{}{}}
\LANGcs{
\Question{
Určete množinu řešení nerovnice.
\[ x^4-1\geq0 \]}
{\( (-\infty;-1\rangle\cup\langle1;\infty) \)}{\( \langle-1;1\rangle \)}{\( \mathbb{R} \)}{\( \langle 1;\infty ) \)}{}{}}
\LANGsk{
\Question{
Určte množinu riešení nerovnice. 
\[ x^4-1\geq0 \]}
{\( (-\infty;-1\rangle\cup\langle1;\infty) \)}{\( \langle-1;1\rangle \)}{\( \mathbb{R} \)}{\( \langle 1;\infty ) \)}{}{}}
\LANGes{
\Question{
Calcula el conjunto de soluciones de la inecuación.
\[ x^4-1\geq0 \]}
{\( (-\infty;-1]\cup[1;\infty) \)}{\( [-1;1] \)}{\( \mathbb{R} \)}{\( [ 1;\infty ) \)}{}{}}
\End

\Data\ID{1003029006}
\Updated{2020-07-15 03:07:12}
\Level{B}
\SubArea{Higher degree equations and inequalities}
\LANGen{
\Question{
Find the solution set of the inequality.
\[ \left(x^4+1\right)\left(x^2+4\right) < 0 \]}
{\( \emptyset \)}{\( (-\infty;-2)\cup(2;\infty) \)}{\( (-2;-1)\cup(1;2) \)}{\( (-1;1) \)}{}{}}
\LANGpl{
\Question{
Wyznacz zbiór rozwiązań podanej nierówności.
\[ \left(x^4+1\right)\left(x^2+4\right) < 0 \]}
{\( \emptyset \)}{\( (-\infty;-2)\cup(2;\infty) \)}{\( (-2;-1)\cup(1;2) \)}{\( (-1;1) \)}{}{}}
\LANGcs{
\Question{
Určete množinu řešení nerovnice.
\[ \left(x^4+1\right)\left(x^2+4\right) < 0 \]}
{\( \emptyset \)}{\( (-\infty;-2)\cup(2;\infty) \)}{\( (-2;-1)\cup(1;2) \)}{\( (-1;1) \)}{}{}}
\LANGsk{
\Question{
Určte množinu riešení nerovnice.
\[ \left(x^4+1\right)\left(x^2+4\right) < 0 \]}
{\( \emptyset \)}{\( (-\infty;-2)\cup(2;\infty) \)}{\( (-2;-1)\cup(1;2) \)}{\( (-1;1) \)}{}{}}
\LANGes{
\Question{
Calcula el conjunto de soluciones de la inecuación.
\[ \left(x^4+1\right)\left(x^2+4\right) < 0 \]}
{\( \emptyset \)}{\( (-\infty;-2)\cup(2;\infty) \)}{\( (-2;-1)\cup(1;2) \)}{\( (-1;1) \)}{}{}}
\End

\Data\ID{1003189101}
\Updated{2020-07-15 03:07:21}
\Level{B}
\SubArea{Higher degree equations and inequalities}
\LANGen{
\Question{
Find the sum of all the positive real roots of the given equation.
\[ x^4-116x^2+1600=0 \]}
{\( 14 \)}{\( 116 \)}{\( 0 \)}{\( 16 \)}{}{}}
\LANGpl{
\Question{
Wyznacz sumę wszystkich rzeczywistych dodatnich  pierwiastków podanego równania. 
\[ x^4-116x^2+1600=0 \]}
{\( 14 \)}{\( 116 \)}{\( 0 \)}{\( 16 \)}{}{}}
\LANGcs{
\Question{
Určete součet všech kladných reálných kořenů dané rovnice. 
\[ x^4-116x^2+1600=0 \]}
{\( 14 \)}{\( 116 \)}{\( 0 \)}{\( 16 \)}{}{}}
\LANGsk{
\Question{
Určte súčet všetkých kladných reálnych koreňov danej rovnice.
\[ x^4-116x^2+1600=0 \]}
{\( 14 \)}{\( 116 \)}{\( 0 \)}{\( 16 \)}{}{}}
\LANGes{
\Question{
Calcula la suma de todas las raíces reales positivas de la siguiente ecuación.
\[ x^4-116x^2+1600=0 \]}
{\( 14 \)}{\( 116 \)}{\( 0 \)}{\( 16 \)}{}{}}
\End

\Data\ID{1003189102}
\Updated{2020-07-15 03:07:21}
\Level{B}
\SubArea{Higher degree equations and inequalities}
\LANGen{
\Question{
Find the product of all the negative real roots of the given equation.
\[ x^4-34x^2+225=0 \]}
{\( 15 \)}{The equation has no negative real roots.}{\( 225 \)}{\( -34 \)}{}{}}
\LANGpl{
\Question{
Wyznacz iloczyn wszystkich rzeczywistych ujemnych pierwiastków podanego równania. 
\[ x^4-34x^2+225=0 \]}
{\( 15 \)}{Równanie nie ma rzeczywistych ujemnych pierwiastków.}{\( 225 \)}{\( -34 \)}{}{}}
\LANGcs{
\Question{
Určete součin všech záporných reálných kořenů dané rovnice.
\[ x^4-34x^2+225=0 \]}
{\( 15 \)}{Rovnice nemá záporné kořeny.}{\( 225 \)}{\( -34 \)}{}{}}
\LANGsk{
\Question{
Určte súčin všetkých záporných reálnych koreňov danej rovnice.
\[ x^4-34x^2+225=0 \]}
{\( 15 \)}{Rovnica nemá záporné reálne korene.}{\( 225 \)}{\( -34 \)}{}{}}
\LANGes{
\Question{
Calcula el producto de todas las raíces reales negativas de la siguiente ecuación.
\[ x^4-34x^2+225=0 \]}
{\( 15 \)}{The equation has no negative real roots.}{\( 225 \)}{\( -34 \)}{}{}}
\End

\Data\ID{1003189103}
\Updated{2020-07-15 03:07:21}
\Level{B}
\SubArea{Higher degree equations and inequalities}
\LANGen{
\Question{
Find the product of all the nonnegative roots of the given equation.
\[ x^4- 13 x^2+ 36=0 \]}
{\( 6 \)}{\( 0 \)}{\( 5 \)}{\( 13 \)}{}{}}
\LANGpl{
\Question{
Wyznacz iloczyn wszystkich nieujemnych pierwiastków podanego równania. 
\[ x^4- 13 x^2+ 36=0 \]}
{\( 6 \)}{\( 0 \)}{\( 5 \)}{\( 13 \)}{}{}}
\LANGcs{
\Question{
Určete součin všech nezáporných reálných kořenů dané rovnice.
\[ x^4- 13 x^2+ 36=0 \]}
{\( 6 \)}{\( 0 \)}{\( 5 \)}{\( 13 \)}{}{}}
\LANGsk{
\Question{
Určte súčin všetkých nezáporných reálnych koreňov danej rovnice.
\[ x^4- 13 x^2+ 36=0 \]}
{\( 6 \)}{\( 0 \)}{\( 5 \)}{\( 13 \)}{}{}}
\LANGes{
\Question{
Calcula el producto de todas las raíces no negativas de la siguiente ecuación.
\[ x^4- 13 x^2+ 36=0 \]}
{\( 6 \)}{\( 0 \)}{\( 5 \)}{\( 13 \)}{}{}}
\End

\Data\ID{9000019803}
\Updated{2020-07-15 03:07:34}
\Level{B}
\SubArea{Higher degree equations and inequalities}
\LANGen{
\Question{
Find the solution set of the following equation. 
\[ 
 x^{4} - 5x^{2} + 4 = 0
\]}
{\(\left \{-2;-1;1;2\right \}\)}{\(\left \{-1;1\right \}\)}{\(\left \{-2;2\right \}\)}{\(\left \{1;2\right \}\)}{}{}}
\LANGpl{
\Question{
Wyznacz zbiór rozwiązań następującego równania.
\[ 
 x^{4} - 5x^{2} + 4 = 0
\]}
{\(\left \{-2;-1;1;2\right \}\)}{\(\left \{-1;1\right \}\)}{\(\left \{-2;2\right \}\)}{\(\left \{1;2\right \}\)}{}{}}
\LANGcs{
\Question{
Určete množinu všech řešení rovnice.
\[x^{4} - 5x^{2} + 4 = 0\]}
{\(\left \{-2;-1;1;2\right \}\)}{\(\left \{-1;1\right \}\)}{\(\left \{-2;2\right \}\)}{\(\left \{1;2\right \}\)}{}{}}
\LANGsk{
\Question{
Určte množinu všetkých riešení danej rovnice. 
\[ 
 x^{4} - 5x^{2} + 4 = 0
\]}
{\(\left \{-2;-1;1;2\right \}\)}{\(\left \{-1;1\right \}\)}{\(\left \{-2;2\right \}\)}{\(\left \{1;2\right \}\)}{}{}}
\LANGes{
\Question{
Calcula el conjunto de soluciones de la siguiente ecuación. 
\[ 
 x^{4} - 5x^{2} + 4 = 0
\]}
{\(\left \{-2;-1;1;2\right \}\)}{\(\left \{-1;1\right \}\)}{\(\left \{-2;2\right \}\)}{\(\left \{1;2\right \}\)}{}{}}
\End

\Data\ID{9000019804}
\Updated{2020-07-15 03:07:34}
\Level{B}
\SubArea{Higher degree equations and inequalities}
\LANGen{
\Question{
Assuming \(x\in \mathbb{R}\), find the solution set of the following equation. 
\[ 
 x^{4} - 16 = 0
\]}
{\(\left \{-2;2\right \}\)}{\(\left \{-\sqrt{2};\sqrt{2}\right \}\)}{\(\left \{-4;4\right \}\)}{\(\left \{-2;-\sqrt{2};\sqrt{2};2\right \}\)}{}{}}
\LANGpl{
\Question{
Zakładając, że \(x\in \mathbb{R}\), wyznacz zbiór rozwiązań następującego równania.
\[ 
 x^{4} - 16 = 0
\]}
{\(\left \{-2;2\right \}\)}{\(\left \{-\sqrt{2};\sqrt{2}\right \}\)}{\(\left \{-4;4\right \}\)}{\(\left \{-2;-\sqrt{2};\sqrt{2};2\right \}\)}{}{}}
\LANGcs{
\Question{
Určete množinu všech řešení rovnice v oboru reálných čísel.
\[x^{4} - 16 = 0\]}
{\(\left \{-2;2\right \}\)}{\(\left \{-\sqrt{2};\sqrt{2}\right \}\)}{\(\left \{-4;4\right \}\)}{\(\left \{-2;-\sqrt{2};\sqrt{2};2\right \}\)}{}{}}
\LANGsk{
\Question{
Určte množinu všetkých riešení danej rovnice pre \(x\in\mathbb{R}\). 
\[ 
 x^{4} - 16 = 0
\]}
{\(\left \{-2;2\right \}\)}{\(\left \{-\sqrt{2};\sqrt{2}\right \}\)}{\(\left \{-4;4\right \}\)}{\(\left \{-2;-\sqrt{2};\sqrt{2};2\right \}\)}{}{}}
\LANGes{
\Question{
Suponiendo que \(x\in \mathbb{R}\), calcula el conjunto de soluciones de la siguiente ecuación. 
\[ 
 x^{4} - 16 = 0
\]}
{\(\left \{-2;2\right \}\)}{\(\left \{-\sqrt{2};\sqrt{2}\right \}\)}{\(\left \{-4;4\right \}\)}{\(\left \{-2;-\sqrt{2};\sqrt{2};2\right \}\)}{}{}}
\End

\Data\ID{9000019805}
\Updated{2020-07-15 03:07:34}
\Level{B}
\SubArea{Higher degree equations and inequalities}
\LANGen{
\Question{
Assuming \(x\in \mathbb{R}\), find the solution set of the following equation. 
\[ 
 x^{4} + 2x^{2} + 1 = 0
\]}
{\(\emptyset \)}{\(\left \{-1;1\right \}\)}{\(\left \{-2;2\right \}\)}{\(\left \{0\right \}\)}{}{}}
\LANGpl{
\Question{
Zakładając, że \(x\in \mathbb{R}\), wyznacz zbiór rozwiązań następującego równania.
\[ 
 x^{4} + 2x^{2} + 1 = 0
\]}
{\(\emptyset \)}{\(\left \{-1;1\right \}\)}{\(\left \{-2;2\right \}\)}{\(\left \{0\right \}\)}{}{}}
\LANGcs{
\Question{
Určete množinu všech řešení rovnice v oboru reálných čísel.
\[x^{4} + 2x^{2} + 1 = 0\]}
{\(\emptyset \)}{\(\left \{-1;1\right \}\)}{\(\left \{-2;2\right \}\)}{\(\left \{0\right \}\)}{}{}}
\LANGsk{
\Question{
Určte množinu všetkých riešení danej rovnice pre \(x\in\mathbb{R}\). 
\[ 
 x^{4} + 2x^{2} + 1 = 0
\]}
{\(\emptyset \)}{\(\left \{-1;1\right \}\)}{\(\left \{-2;2\right \}\)}{\(\left \{0\right \}\)}{}{}}
\LANGes{
\Question{
Suponiendo que \(x\in \mathbb{R}\), calcula el conjunto de soluciones de la siguiente ecuación.
\[ 
 x^{4} + 2x^{2} + 1 = 0
\]}
{\(\emptyset \)}{\(\left \{-1;1\right \}\)}{\(\left \{-2;2\right \}\)}{\(\left \{0\right \}\)}{}{}}
\End

\Data\ID{9000019806}
\Updated{2020-07-15 03:07:34}
\Level{B}
\SubArea{Higher degree equations and inequalities}
\LANGen{
\Question{
Find the smallest integer solution of the following equation. 
\[ 
 x^{4} - 2x^{3} - x^{2} + 2x = 0
\]}
{\(- 1\)}{\(0\)}{\(1\)}{\(2\)}{}{}}
\LANGpl{
\Question{
Znajdź najmniejszą liczbę całkowitą, która jest rozwiązaniem podanego równania.
\[ 
 x^{4} - 2x^{3} - x^{2} + 2x = 0
\]}
{\(- 1\)}{\(0\)}{\(1\)}{\(2\)}{}{}}
\LANGcs{
\Question{
Najděte nejmenší číslo \(x\in \mathbb{Z}\),
které je řešením rovnice \(x^{4} - 2x^{3} - x^{2} + 2x = 0\).}
{\(- 1\)}{\(0\)}{\(1\)}{\(2\)}{}{}}
\LANGsk{
\Question{
Nájdite najmenšie celé číslo, ktoré je riešením danej rovnice. 
\[ 
 x^{4} - 2x^{3} - x^{2} + 2x = 0
\]}
{\(- 1\)}{\(0\)}{\(1\)}{\(2\)}{}{}}
\LANGes{
\Question{
Halla el número más pequeño \(x\in \mathbb{Z}\),
que sea solución de la siguiente ecuación \(x^{4} - 2x^{3} - x^{2} + 2x = 0\).}
{\(- 1\)}{\(0\)}{\(1\)}{\(2\)}{}{}}
\End

\Data\ID{9000019809}
\Updated{2020-07-15 03:07:34}
\Level{B}
\SubArea{Higher degree equations and inequalities}
\LANGen{
\Question{
Find the factorization of the following equation. 
\[ 
 x^{3} + 3x^{2} - x - 3 = 0
\]}
{\(\left (x + 3\right )\left (x + 1\right )\left (x - 1\right ) = 0\)}{\(\left (x - 3\right )\left (x + 1\right )\left (x - 1\right ) = 0\)}{\(\left (x + 3\right )\left (x - 3\right )\left (x - 1\right ) = 0\)}{\(\left (x + 3\right )\left (x - 3\right )\left (x + 1\right ) = 0\)}{}{}}
\LANGpl{
\Question{
Rozłóż na czynniki pierwsze podane równanie.
\[ 
 x^{3} + 3x^{2} - x - 3 = 0
\]}
{\(\left (x + 3\right )\left (x + 1\right )\left (x - 1\right ) = 0\)}{\(\left (x - 3\right )\left (x + 1\right )\left (x - 1\right ) = 0\)}{\(\left (x + 3\right )\left (x - 3\right )\left (x - 1\right ) = 0\)}{\(\left (x + 3\right )\left (x - 3\right )\left (x + 1\right ) = 0\)}{}{}}
\LANGcs{
\Question{
Vyberte správný součinový tvar dané rovnice.
\[x^{3} + 3x^{2} - x - 3 = 0\]}
{\(\left (x + 3\right )\left (x + 1\right )\left (x - 1\right ) = 0\)}{\(\left (x - 3\right )\left (x + 1\right )\left (x - 1\right ) = 0\)}{\(\left (x + 3\right )\left (x - 3\right )\left (x - 1\right ) = 0\)}{\(\left (x + 3\right )\left (x - 3\right )\left (x + 1\right ) = 0\)}{}{}}
\LANGsk{
\Question{
Vyberte správny súčinový tvar danej rovnice.
\[ 
 x^{3} + 3x^{2} - x - 3 = 0
\]}
{\(\left (x + 3\right )\left (x + 1\right )\left (x - 1\right ) = 0\)}{\(\left (x - 3\right )\left (x + 1\right )\left (x - 1\right ) = 0\)}{\(\left (x + 3\right )\left (x - 3\right )\left (x - 1\right ) = 0\)}{\(\left (x + 3\right )\left (x - 3\right )\left (x + 1\right ) = 0\)}{}{}}
\LANGes{
\Question{
Halla la factorización de la siguiente ecuación. 
\[ 
 x^{3} + 3x^{2} - x - 3 = 0
\]}
{\(\left (x + 3\right )\left (x + 1\right )\left (x - 1\right ) = 0\)}{\(\left (x - 3\right )\left (x + 1\right )\left (x - 1\right ) = 0\)}{\(\left (x + 3\right )\left (x - 3\right )\left (x - 1\right ) = 0\)}{\(\left (x + 3\right )\left (x - 3\right )\left (x + 1\right ) = 0\)}{}{}}
\End

\Data\ID{9000019810}
\Updated{2020-07-15 03:07:34}
\Level{B}
\SubArea{Higher degree equations and inequalities}
\LANGen{
\Question{
Find the factorization of the following equation. 
\[ 
 5x^{4} - 30x^{2} + 40 = 0
\]}
{\(5\left (x -\sqrt{2}\right )\left (x + \sqrt{2}\right )\left (x - 2\right )\left (x + 2\right ) = 0\)}{\(\left (x -\sqrt{2}\right )\left (x + \sqrt{2}\right )\left (x - 2\right )\left (x + 2\right ) = 0\)}{\(5x\left (x -\sqrt{2}\right )\left (x + \sqrt{2}\right )\left (x - 2\right ) = 0\)}{\(5x\left (x -\sqrt{2}\right )\left (x + \sqrt{2}\right )\left (x + 2\right ) = 0\)}{}{}}
\LANGpl{
\Question{
Rozłóż na czynniki pierwsze podane równanie.
\[ 
 5x^{4} - 30x^{2} + 40 = 0
\]}
{\(5\left (x -\sqrt{2}\right )\left (x + \sqrt{2}\right )\left (x - 2\right )\left (x + 2\right ) = 0\)}{\(\left (x -\sqrt{2}\right )\left (x + \sqrt{2}\right )\left (x - 2\right )\left (x + 2\right ) = 0\)}{\(5x\left (x -\sqrt{2}\right )\left (x + \sqrt{2}\right )\left (x - 2\right ) = 0\)}{\(5x\left (x -\sqrt{2}\right )\left (x + \sqrt{2}\right )\left (x + 2\right ) = 0\)}{}{}}
\LANGcs{
\Question{
Vyberte správný součinový tvar dané rovnice.
\[5x^{4} - 30x^{2} + 40 = 0\]}
{\(5\left (x -\sqrt{2}\right )\left (x + \sqrt{2}\right )\left (x - 2\right )\left (x + 2\right ) = 0\)}{\(\left (x -\sqrt{2}\right )\left (x + \sqrt{2}\right )\left (x - 2\right )\left (x + 2\right ) = 0\)}{\(5x\left (x -\sqrt{2}\right )\left (x + \sqrt{2}\right )\left (x - 2\right ) = 0\)}{\(5x\left (x -\sqrt{2}\right )\left (x + \sqrt{2}\right )\left (x + 2\right ) = 0\)}{}{}}
\LANGsk{
\Question{
Vyberte správny súčinový tvar danej rovnice. 
\[ 
 5x^{4} - 30x^{2} + 40 = 0
\]}
{\(5\left (x -\sqrt{2}\right )\left (x + \sqrt{2}\right )\left (x - 2\right )\left (x + 2\right ) = 0\)}{\(\left (x -\sqrt{2}\right )\left (x + \sqrt{2}\right )\left (x - 2\right )\left (x + 2\right ) = 0\)}{\(5x\left (x -\sqrt{2}\right )\left (x + \sqrt{2}\right )\left (x - 2\right ) = 0\)}{\(5x\left (x -\sqrt{2}\right )\left (x + \sqrt{2}\right )\left (x + 2\right ) = 0\)}{}{}}
\LANGes{
\Question{
Halla la factorización de la siguiente ecuación..
\[ 
 5x^{4} - 30x^{2} + 40 = 0
\]}
{\(5\left (x -\sqrt{2}\right )\left (x + \sqrt{2}\right )\left (x - 2\right )\left (x + 2\right ) = 0\)}{\(\left (x -\sqrt{2}\right )\left (x + \sqrt{2}\right )\left (x - 2\right )\left (x + 2\right ) = 0\)}{\(5x\left (x -\sqrt{2}\right )\left (x + \sqrt{2}\right )\left (x - 2\right ) = 0\)}{\(5x\left (x -\sqrt{2}\right )\left (x + \sqrt{2}\right )\left (x + 2\right ) = 0\)}{}{}}
\End

\Data\ID{9000028301}
\Updated{2020-07-15 03:07:34}
\Level{B}
\SubArea{Higher degree equations and inequalities}
\LANGen{
\Question{
The following equation has a solution
\(x = 1\). Find
the sum of the remaining real solutions. 
\[ 
 x^{3} - 7x + 6 = 0
\]}
{\(- 1\)}{\(1\)}{\(0\)}{\(2\)}{}{}}
\LANGpl{
\Question{
Rozwiązaniem podanego równania jest 
\(x = 1\). Oblicz sumę pozostałych rzeczywistych rozwiązań.
\[ 
 x^{3} - 7x + 6 = 0
\]}
{\(- 1\)}{\(1\)}{\(0\)}{\(2\)}{}{}}
\LANGcs{
\Question{
Jeden kořen následující rovnice je roven $1$. Určete součet zbývajících reálných kořenů.
\[x^{3} - 7x + 6 = 0\]}
{\(- 1\)}{\(1\)}{\(0\)}{\(2\)}{}{}}
\LANGsk{
\Question{
Nasledujúca rovnica má jeden koreň
\(x = 1\). Určte súčet zvyšných reálnych koreňov.
\[ 
 x^{3} - 7x + 6 = 0
\]}
{\(- 1\)}{\(1\)}{\(0\)}{\(2\)}{}{}}
\LANGes{
\Question{
La siguiente ecuación tiene una solución \(x = 1\). Calcula la suma de las soluciones reales restantes. 
\[ 
 x^{3} - 7x + 6 = 0
\]}
{\(- 1\)}{\(1\)}{\(0\)}{\(2\)}{}{}}
\End

\Data\ID{9000028302}
\Updated{2020-07-15 03:07:34}
\Level{B}
\SubArea{Higher degree equations and inequalities}
\LANGen{
\Question{
The following equation has a solution
\(x = 1\). Find
the sum of the remaining real solutions. 
\[ 
 x^{3} + 2x^{2} - x - 2 = 0
\]}
{\(- 3\)}{\(- 1\)}{\(0\)}{\(2\)}{}{}}
\LANGpl{
\Question{
Rozwiązaniem podanego równania jest 
\(x = 1\). Oblicz sumę pozostałych rzeczywistych rozwiązań.
\[ 
 x^{3} + 2x^{2} - x - 2 = 0
\]}
{\(- 3\)}{\(- 1\)}{\(0\)}{\(2\)}{}{}}
\LANGcs{
\Question{
Jeden z kořen následující rovnice je roven $1$. Určete součet zbývajících reálných kořenů.
\[x^{3} + 2x^{2} - x - 2 = 0\]}
{\(- 3\)}{\(- 1\)}{\(0\)}{\(2\)}{}{}}
\LANGsk{
\Question{
Daná rovnica má jeden koreň
\(x = 1\). Určte
súčet zvyšných reálnych koreňov rovnice.  
\[ 
 x^{3} + 2x^{2} - x - 2 = 0
\]}
{\(- 3\)}{\(- 1\)}{\(0\)}{\(2\)}{}{}}
\LANGes{
\Question{
La siguiente ecuación tiene una solución \(x = 1\). Calcula la suma de las soluciones reales restantes. 
\[ 
 x^{3} + 2x^{2} - x - 2 = 0
\]}
{\(- 3\)}{\(- 1\)}{\(0\)}{\(2\)}{}{}}
\End

\Data\ID{9000028303}
\Updated{2020-07-15 03:07:34}
\Level{B}
\SubArea{Higher degree equations and inequalities}
\LANGen{
\Question{
The following equation has a solution
\(x = -2\). Find
the sum of the remaining real solutions. 
\[ 
 x^{3} + 3x^{2} - 18x - 40 = 0
\]}
{\(- 1\)}{\(1\)}{\(0\)}{\(4\)}{}{}}
\LANGpl{
\Question{
Rozwiązaniem podanego równania jest 
\(x = -2\).  Oblicz sumę pozostałych rzeczywistych rozwiązań.
\[ 
 x^{3} + 3x^{2} - 18x - 40 = 0
\]}
{\(- 1\)}{\(1\)}{\(0\)}{\(4\)}{}{}}
\LANGcs{
\Question{
Jeden kořen následující rovnice je roven $-2$. Určete součet zbývajících reálných kořenů.
\[x^{3} + 3x^{2} - 18x - 40 = 0\]}
{\(- 1\)}{\(1\)}{\(0\)}{\(4\)}{}{}}
\LANGsk{
\Question{
Daná rovnica má jeden koreň
\(x = -2\). Určte súčet 
zvyšných reálnych koreňov rovnice.
\[ 
 x^{3} + 3x^{2} - 18x - 40 = 0
\]}
{\(- 1\)}{\(1\)}{\(0\)}{\(4\)}{}{}}
\LANGes{
\Question{
La siguiente ecuación tiene una solución \(x = -2\). Calcula la suma de las soluciones reales restantes.
\[ 
 x^{3} + 3x^{2} - 18x - 40 = 0
\]}
{\(- 1\)}{\(1\)}{\(0\)}{\(4\)}{}{}}
\End

\Data\ID{9000028307}
\Updated{2020-07-15 03:07:34}
\Level{B}
\SubArea{Higher degree equations and inequalities}
\LANGen{
\Question{
Solve the following equation.
\[ 
 x^{3} + 6x^{2} - 8x = 0
\]}
{\(0\),
\(- 3 -\sqrt{17}\),
\(- 3 + \sqrt{17}\)}{\(0\),
\(3 -\sqrt{17}\),
\(3 + \sqrt{17}\)}{\(0\),
\(- 3\),
\(\sqrt{
17}\)}{\(0\),
\(3\),
\(-\sqrt{17}\)}{}{}}
\LANGpl{
\Question{
Rozwiąż podane równanie.
\[ 
 x^{3} + 6x^{2} - 8x = 0
\]}
{\(0\),
\(- 3 -\sqrt{17}\),
\(- 3 + \sqrt{17}\)}{\(0\),
\(3 -\sqrt{17}\),
\(3 + \sqrt{17}\)}{\(0\),
\(- 3\),
\(\sqrt{
17}\)}{\(0\),
\(3\),
\(-\sqrt{17}\)}{}{}}
\LANGcs{
\Question{
Najděte všechna řešení dané rovnice.
\[ 
 x^{3} + 6x^{2} - 8x = 0
\]}
{\(0\),
\(- 3 -\sqrt{17}\),
\(- 3 + \sqrt{17}\)}{\(0\),
\(3 -\sqrt{17}\),
\(3 + \sqrt{17}\)}{\(0\),
\(- 3\),
\(\sqrt{
17}\)}{\(0\),
\(3\),
\(-\sqrt{17}\)}{}{}}
\LANGsk{
\Question{
Vyriešte danú rovnicu.
\[ 
 x^{3} + 6x^{2} - 8x = 0
\]}
{\(0\),
\(- 3 -\sqrt{17}\),
\(- 3 + \sqrt{17}\)}{\(0\),
\(3 -\sqrt{17}\),
\(3 + \sqrt{17}\)}{\(0\),
\(- 3\),
\(\sqrt{
17}\)}{\(0\),
\(3\),
\(-\sqrt{17}\)}{}{}}
\LANGes{
\Question{
Resuelve la siguiente ecuación.
\[ 
 x^{3} + 6x^{2} - 8x = 0
\]}
{\(0\),
\(- 3 -\sqrt{17}\),
\(- 3 + \sqrt{17}\)}{\(0\),
\(3 -\sqrt{17}\),
\(3 + \sqrt{17}\)}{\(0\),
\(- 3\),
\(\sqrt{
17}\)}{\(0\),
\(3\),
\(-\sqrt{17}\)}{}{}}
\End

\Data\ID{9000028308}
\Updated{2020-07-15 03:07:34}
\Level{B}
\SubArea{Higher degree equations and inequalities}
\LANGen{
\Question{
Solve the following equation.
\[ 
 x^{4} - 20x^{2} + 99 = 0
\]}
{\(-\sqrt{11}\),
\(- 3\),
\(3\),
\(\sqrt{
11}\)}{\(0\),
\(- 3 -\sqrt{17}\),
\(- 3 + \sqrt{17}\)}{\(0\),
\(3 -\sqrt{17}\),
\(3 + \sqrt{17}\)}{\(-\sqrt{17}\),
\(- 3\),
\(3\),
\(\sqrt{
17}\)}{}{}}
\LANGpl{
\Question{
Rozwiąż podane równanie.
\[ 
 x^{4} - 20x^{2} + 99 = 0
\]}
{\(-\sqrt{11}\),
\(- 3\),
\(3\),
\(\sqrt{
11}\)}{\(0\),
\(- 3 -\sqrt{17}\),
\(- 3 + \sqrt{17}\)}{\(0\),
\(3 -\sqrt{17}\),
\(3 + \sqrt{17}\)}{\(-\sqrt{17}\),
\(- 3\),
\(3\),
\(\sqrt{
17}\)}{}{}}
\LANGcs{
\Question{
Najděte všechna řešení dané rovnice.
\[ 
 x^{4} - 20x^{2} + 99 = 0
\]}
{\(-\sqrt{11}\),
\(- 3\),
\(3\),
\(\sqrt{
11}\)}{\(0\),
\(- 3 -\sqrt{17}\),
\(- 3 + \sqrt{17}\)}{\(0\),
\(3 -\sqrt{17}\),
\(3 + \sqrt{17}\)}{\(-\sqrt{17}\),
\(- 3\),
\(3\),
\(\sqrt{
17}\)}{}{}}
\LANGsk{
\Question{
Vyriešte danú rovnicu.
\[ 
 x^{4} - 20x^{2} + 99 = 0
\]}
{\(-\sqrt{11}\),
\(- 3\),
\(3\),
\(\sqrt{
11}\)}{\(0\),
\(- 3 -\sqrt{17}\),
\(- 3 + \sqrt{17}\)}{\(0\),
\(3 -\sqrt{17}\),
\(3 + \sqrt{17}\)}{\(-\sqrt{17}\),
\(- 3\),
\(3\),
\(\sqrt{
17}\)}{}{}}
\LANGes{
\Question{
Resuelve la siguiente ecuación.
\[ 
 x^{4} - 20x^{2} + 99 = 0
\]}
{\(-\sqrt{11}\),
\(- 3\),
\(3\),
\(\sqrt{
11}\)}{\(0\),
\(- 3 -\sqrt{17}\),
\(- 3 + \sqrt{17}\)}{\(0\),
\(3 -\sqrt{17}\),
\(3 + \sqrt{17}\)}{\(-\sqrt{17}\),
\(- 3\),
\(3\),
\(\sqrt{
17}\)}{}{}}
\End

\Data\ID{9000028310}
\Updated{2020-07-15 03:07:34}
\Level{B}
\SubArea{Higher degree equations and inequalities}
\LANGen{
\Question{
Find the sum of all real solutions of the following equation. 
\[ 
 x^{4} - 20x^{2} + 64 = 0
\]}
{\(0\)}{\(- 10\)}{\(4\)}{\(10\)}{}{}}
\LANGpl{
\Question{
Oblicz sumę wszystkich rzeczywistych rozwiązań podanego równania.
\[ 
 x^{4} - 20x^{2} + 64 = 0
\]}
{\(0\)}{\(- 10\)}{\(4\)}{\(10\)}{}{}}
\LANGcs{
\Question{
Určete součet všech reálných kořenů dané rovnice.
\[ 
 x^{4} - 20x^{2} + 64 = 0
\]}
{\(0\)}{\(- 10\)}{\(4\)}{\(10\)}{}{}}
\LANGsk{
\Question{
Určte súčet všetkých reálnych koreňov danej rovnice. 
\[ 
 x^{4} - 20x^{2} + 64 = 0
\]}
{\(0\)}{\(- 10\)}{\(4\)}{\(10\)}{}{}}
\LANGes{
\Question{
Calcula la suma de todas las soluciones reales de la siguiente ecuación.
\[ 
 x^{4} - 20x^{2} + 64 = 0
\]}
{\(0\)}{\(- 10\)}{\(4\)}{\(10\)}{}{}}
\End

\Data\ID{9000029301}
\Updated{2020-07-15 03:07:34}
\Level{B}
\SubArea{Higher degree equations and inequalities}
\LANGen{
\Question{
Find the solution set of the following inequality. 
\[ 
 \left (x - 1\right )\left (x - 2\right )\left (x - 3\right )\geq 0
\]}
{\(\left [ 1;2\right ] \cup \left [ 3;\infty \right )\)}{\(\left (-\infty ;\infty \right )\)}{\(\left (-\infty ;1\right )\cup \left (2;3\right )\)}{\(\emptyset \)}{\(\{0\}\)}{}}
\LANGpl{
\Question{
Wyznacz zbiór rozwiązań podanej nierówności. 
\[ 
 \left (x - 1\right )\left (x - 2\right )\left (x - 3\right )\geq 0
\]}
{\(\left [ 1;2\right ] \cup \left [ 3;\infty \right )\)}{\(\left (-\infty ;\infty \right )\)}{\(\left (-\infty ;1\right )\cup \left (2;3\right )\)}{\(\emptyset \)}{\(\{0\}\)}{}}
\LANGcs{
\Question{
Vyberte řešení dané nerovnice. \[\left (x - 1\right )\left (x - 2\right )\left (x - 3\right )\geq 0\]}
{\(\left \langle 1;2\right \rangle \cup \left \langle 3;\infty \right )\)}{\(\left (-\infty ;\infty \right )\)}{\(\left (-\infty ;1\right )\cup \left (2;3\right )\)}{\(\emptyset \)}{\(\{0\}\)}{}}
\LANGsk{
\Question{
Vyberte riešenie danej nerovnice. 
\[ 
 \left (x - 1\right )\left (x - 2\right )\left (x - 3\right )\geq 0
\]}
{\(\left [ 1;2\right ] \cup \left [ 3;\infty \right )\)}{\(\left (-\infty ;\infty \right )\)}{\(\left (-\infty ;1\right )\cup \left (2;3\right )\)}{\(\emptyset \)}{\(\{0\}\)}{}}
\LANGes{
\Question{
Calcula el conjunto de soluciones de la siguiente inecuación. 
\[ 
 \left (x - 1\right )\left (x - 2\right )\left (x - 3\right )\geq 0
\]}
{\(\left [ 1;2\right ] \cup \left [ 3;\infty \right )\)}{\(\left (-\infty ;\infty \right )\)}{\(\left (-\infty ;1\right )\cup \left (2;3\right )\)}{\(\emptyset \)}{\(\{0\}\)}{}}
\End

\Data\ID{9000029302}
\Updated{2020-07-15 03:07:34}
\Level{B}
\SubArea{Higher degree equations and inequalities}
\LANGen{
\Question{
Find the solution set of the following inequality. 
\[ 
 x^{4} - 16 > 0
\]}
{\(\mathbb{R}\setminus \left [ -2;2\right ] \)}{\(\mathbb{R}\)}{\(\left (-\infty ;-4\right )\cup \left (4;\infty \right )\)}{\(\left (-2;2\right )\)}{\(\left (-4;4\right )\)}{}}
\LANGpl{
\Question{
Wyznacz zbiór rozwiązań podanej nierówności. 
\[ 
 x^{4} - 16 > 0
\]}
{\(\mathbb{R}\setminus \left [ -2;2\right ] \)}{\(\mathbb{R}\)}{\(\left (-\infty ;-4\right )\cup \left (4;\infty \right )\)}{\(\left (-2;2\right )\)}{\(\left (-4;4\right )\)}{}}
\LANGcs{
\Question{
Najděte řešení nerovnice \(x^{4} - 16 > 0\).}
{\(\mathbb{R}\setminus \left \langle -2;2\right \rangle \)}{\(\mathbb{R}\)}{\(\left (-\infty ;-4\right )\cup \left (4;\infty \right )\)}{\(\left (-2;2\right )\)}{\(\left (-4;4\right )\)}{}}
\LANGsk{
\Question{
Nájdite riešenie danej nerovnice. 
\[ 
 x^{4} - 16 > 0
\]}
{\(\mathbb{R}\setminus \left [ -2;2\right ] \)}{\(\mathbb{R}\)}{\(\left (-\infty ;-4\right )\cup \left (4;\infty \right )\)}{\(\left (-2;2\right )\)}{\(\left (-4;4\right )\)}{}}
\LANGes{
\Question{
Calcula el conjunto de soluciones de la siguiente inecuación. 
\[ 
 x^{4} - 16 > 0
\]}
{\(\mathbb{R}\setminus \left [ -2;2\right ] \)}{\(\mathbb{R}\)}{\(\left (-\infty ;-4\right )\cup \left (4;\infty \right )\)}{\(\left (-2;2\right )\)}{\(\left (-4;4\right )\)}{}}
\End

\Data\ID{9000029303}
\Updated{2020-07-15 03:07:34}
\Level{B}
\SubArea{Higher degree equations and inequalities}
\LANGen{
\Question{
In the following list identify an inequality which does not have solution in
\(\mathbb{R}\).}
{\(x^{4} + 81 < 0\)}{\((x - 3)^{3} > 0\)}{\(x^{3} - 9x < 0\)}{\(4x^{4} - 64 > 0\)}{}{}}
\LANGpl{
\Question{
Która z poniższych nierówności nie ma rozwiązania w 
\(\mathbb{R}\)?}
{\(x^{4} + 81 < 0\)}{\((x - 3)^{3} > 0\)}{\(x^{3} - 9x < 0\)}{\(4x^{4} - 64 > 0\)}{}{}}
\LANGcs{
\Question{
Vyberte z uvedených nerovnic tu, která nemá řešení v oboru reálných
čísel.}
{\(x^{4} + 81 < 0\)}{\((x - 3)^{3} > 0\)}{\(x^{3} - 9x < 0\)}{\(4x^{4} - 64 > 0\)}{}{}}
\LANGsk{
\Question{
Z uvedených nerovníc vyberte tú, ktorá nemá riešenie v 
\(\mathbb{R}\).}
{\(x^{4} + 81 < 0\)}{\((x - 3)^{3} > 0\)}{\(x^{3} - 9x < 0\)}{\(4x^{4} - 64 > 0\)}{}{}}
\LANGes{
\Question{
En la siguiente lista, identifica la inecuación que no tiene solución en
\(\mathbb{R}\).}
{\(x^{4} + 81 < 0\)}{\((x - 3)^{3} > 0\)}{\(x^{3} - 9x < 0\)}{\(4x^{4} - 64 > 0\)}{}{}}
\End

\Data\ID{9000029304}
\Updated{2020-07-15 03:07:34}
\Level{B}
\SubArea{Higher degree equations and inequalities}
\LANGen{
\Question{
Find the solution set of the following inequality. 
\[ 
 x^{3} - 3x^{2} + 2x\geq 0
\]}
{\(\left [ 0;1\right ] \cup \left [ 2;\infty \right )\)}{\(\mathbb{R}\)}{\(\emptyset \)}{\(\left (-\infty ;0\right ] \cup \left [ 1;2\right ] \)}{}{}}
\LANGpl{
\Question{
Wyznacz zbiór rozwiązań podanej nierówności. 
\[ 
 x^{3} - 3x^{2} + 2x\geq 0
\]}
{\(\left [ 0;1\right ] \cup \left [ 2;\infty \right )\)}{\(\mathbb{R}\)}{\(\emptyset \)}{\(\left (-\infty ;0\right ] \cup \left [ 1;2\right ] \)}{}{}}
\LANGcs{
\Question{
Vyberte množinu řešení následující nerovnice.
\[x^{3} - 3x^{2} + 2x\geq 0\]}
{\(\left \langle 0;1\right \rangle \cup \left \langle 2;\infty \right )\)}{\(\mathbb{R}\)}{$\emptyset$}{\(\left (-\infty ;0\right \rangle \cup \left \langle 1;2\right \rangle \)}{}{}}
\LANGsk{
\Question{
Nájdite riešenie danej nerovnice. 
\[ 
 x^{3} - 3x^{2} + 2x\geq 0
\]}
{\(\left [ 0;1\right ] \cup \left [ 2;\infty \right )\)}{\(\mathbb{R}\)}{\(\emptyset \)}{\(\left (-\infty ;0\right ] \cup \left [ 1;2\right ] \)}{}{}}
\LANGes{
\Question{
Calcula el conjunto de soluciones de la siguiente inecuación. 
\[ 
 x^{3} - 3x^{2} + 2x\geq 0
\]}
{\(\left [ 0;1\right ] \cup \left [ 2;\infty \right )\)}{\(\mathbb{R}\)}{\(\emptyset \)}{\(\left (-\infty ;0\right ] \cup \left [ 1;2\right ] \)}{}{}}
\End

\Data\ID{9000029305}
\Updated{2020-07-15 03:07:34}
\Level{B}
\SubArea{Higher degree equations and inequalities}
\LANGen{
\Question{
Find the solution set of the following inequality. 
\[ 
 x^{4} + 81\leq 0
\]}
{\(\emptyset \)}{\(0\)}{\(\mathbb{R}\setminus \left (-9;9\right )\)}{\(\mathbb{R}\)}{\(\left (-\infty ;-3\right ] \cup \left [ 3;\infty \right )\)}{}}
\LANGpl{
\Question{
Wyznacz zbiór rozwiązań podanej nierówności. 
\[ 
 x^{4} + 81\leq 0
\]}
{\(\emptyset \)}{\(0\)}{\(\mathbb{R}\setminus \left (-9;9\right )\)}{\(\mathbb{R}\)}{\(\left (-\infty ;-3\right ] \cup \left [ 3;\infty \right )\)}{}}
\LANGcs{
\Question{
Najděte řešení nerovnice \(x^{4} + 81\leq 0\).}
{\(\emptyset \)}{\(0\)}{\(\mathbb{R}\setminus \left (-9;9\right )\)}{\(\mathbb{R}\)}{\(\left (-\infty ;-3\right \rangle \cup \left \langle 3;\infty \right )\)}{}}
\LANGsk{
\Question{
Nájdite riešenie danej nerovnice. 
\[ 
 x^{4} + 81\leq 0
\]}
{\(\emptyset \)}{\(0\)}{\(\mathbb{R}\setminus \left (-9;9\right )\)}{\(\mathbb{R}\)}{\(\left (-\infty ;-3\right ] \cup \left [ 3;\infty \right )\)}{}}
\LANGes{
\Question{
Calcula el conjunto de soluciones de la siguiente inecuación. 
\[ 
 x^{4} + 81\leq 0
\]}
{\(\emptyset \)}{\(0\)}{\(\mathbb{R}\setminus \left (-9;9\right )\)}{\(\mathbb{R}\)}{\(\left (-\infty ;-3\right ] \cup \left [ 3;\infty \right )\)}{}}
\End

\Data\ID{9000029307}
\Updated{2020-07-15 03:07:34}
\Level{B}
\SubArea{Higher degree equations and inequalities}
\LANGen{
\Question{
In the following list identify an inequality which holds for every
\(x\in \mathbb{R}\).}
{\(- x^{4} - x^{2}\leq 0\)}{\(x^{3} + 3x^{2} + 3x + 1 > 0\)}{\(x^{4} + x^{2} + 1 < 0\)}{\(- x^{3} + 6x^{2} - 12x + 8 > 0\)}{}{}}
\LANGpl{
\Question{
Rozwiązaniem której z poniższych nierówności jest każde
\(x\in \mathbb{R}\)?}
{\(- x^{4} - x^{2}\leq 0\)}{\(x^{3} + 3x^{2} + 3x + 1 > 0\)}{\(x^{4} + x^{2} + 1 < 0\)}{\(- x^{3} + 6x^{2} - 12x + 8 > 0\)}{}{}}
\LANGcs{
\Question{
Vyberte z uvedených nerovnic tu, jejímž řešením je množina všech
reálných čísel.}
{\(- x^{4} - x^{2}\leq 0\)}{\(x^{3} + 3x^{2} + 3x + 1 > 0\)}{\(x^{4} + x^{2} + 1 < 0\)}{\(- x^{3} + 6x^{2} - 12x + 8 > 0\)}{}{}}
\LANGsk{
\Question{
Vyberte z uvedených nerovníc tú, ktorej riešením je množina všetkých reálnych čísel.}
{\(- x^{4} - x^{2}\leq 0\)}{\(x^{3} + 3x^{2} + 3x + 1 > 0\)}{\(x^{4} + x^{2} + 1 < 0\)}{\(- x^{3} + 6x^{2} - 12x + 8 > 0\)}{}{}}
\LANGes{
\Question{
Identifica la inecuación cuya solución es el conjunto de todos los números reales.}
{\(- x^{4} - x^{2}\leq 0\)}{\(x^{3} + 3x^{2} + 3x + 1 > 0\)}{\(x^{4} + x^{2} + 1 < 0\)}{\(- x^{3} + 6x^{2} - 12x + 8 > 0\)}{}{}}
\End

\Data\ID{9000029308}
\Updated{2020-07-15 03:07:34}
\Level{B}
\SubArea{Higher degree equations and inequalities}
\LANGen{
\Question{
Find the solution set of the following inequality. 
\[ 
 x^{3} + 4x < 0
\]}
{\((-\infty ;0)\)}{\((2;\infty )\)}{\(\mathbb{R}\)}{\(\emptyset \)}{\((0;\infty )\)}{}}
\LANGpl{
\Question{
Wyznacz zbiór rozwiązań podanej nierówności. 
\[ 
 x^{3} + 4x < 0
\]}
{\((-\infty ;0)\)}{\((2;\infty )\)}{\(\mathbb{R}\)}{\(\emptyset \)}{\((0;\infty )\)}{}}
\LANGcs{
\Question{
Vyberte množinu řešení následující nerovnice.
\[x^{3} + 4x < 0\]}
{\((-\infty ;0)\)}{\((2;\infty )\)}{\(\mathbb{R}\)}{$\emptyset$}{\((0;\infty )\)}{}}
\LANGsk{
\Question{
Nájdite riešenie danej nerovnice. 
\[ 
 x^{3} + 4x < 0
\]}
{\((-\infty ;0)\)}{\((2;\infty )\)}{\(\mathbb{R}\)}{\(\emptyset \)}{\((0;\infty )\)}{}}
\LANGes{
\Question{
Calcula el conjunto de soluciones de la siguiente inecuación. 
\[ 
 x^{3} + 4x < 0
\]}
{\((-\infty ;0)\)}{\((2;\infty )\)}{\(\mathbb{R}\)}{\(\emptyset \)}{\((0;\infty )\)}{}}
\End

\Data\ID{9000029309}
\Updated{2020-07-15 03:07:34}
\Level{B}
\SubArea{Higher degree equations and inequalities}
\LANGen{
\Question{
Find the solution set of the following inequality. 
\[ 
 (x - 1)(x - 2)(x - 3) < (x - 1)(x - 2)
\]}
{\((-\infty ;1)\cup (2;4)\)}{\(\emptyset \)}{\((0;3)\)}{\((-\infty ;-3)\cup (3;\infty )\)}{\((-3;3)\)}{}}
\LANGpl{
\Question{
Wyznacz zbiór rozwiązań podanej nierówności. 
\[ 
 (x - 1)(x - 2)(x - 3) < (x - 1)(x - 2)
\]}
{\((-\infty ;1)\cup (2;4)\)}{\(\emptyset \)}{\((0;3)\)}{\((-\infty ;-3)\cup (3;\infty )\)}{\((-3;3)\)}{}}
\LANGcs{
\Question{
Vyberte množinu řešení následující nerovnice.
\[(x - 1)(x - 2)(x - 3) < (x - 1)(x - 2)\]}
{\((-\infty ;1)\cup (2;4)\)}{$\emptyset$}{\((0;3)\)}{\((-\infty ;-3)\cup (3;\infty )\)}{\((-3;3)\)}{}}
\LANGsk{
\Question{
Nájdite riešenie danej nerovnice. 
\[ 
 (x - 1)(x - 2)(x - 3) < (x - 1)(x - 2)
\]}
{\((-\infty ;1)\cup (2;4)\)}{\(\emptyset \)}{\((0;3)\)}{\((-\infty ;-3)\cup (3;\infty )\)}{\((-3;3)\)}{}}
\LANGes{
\Question{
Calcula el conjunto de soluciones de la siguiente inecuación. 
\[ 
 (x - 1)(x - 2)(x - 3) < (x - 1)(x - 2)
\]}
{\((-\infty ;1)\cup (2;4)\)}{\(\emptyset \)}{\((0;3)\)}{\((-\infty ;-3)\cup (3;\infty )\)}{\((-3;3)\)}{}}
\End

\Data\ID{9000029310}
\Updated{2020-07-15 03:07:34}
\Level{B}
\SubArea{Higher degree equations and inequalities}
\LANGen{
\Question{
Find the solution set of the following inequality.
\[ 
 (x + 2)(x^{2} + 4x + 3) > x^{2} + 5x + 6
\]}
{\((-3;-2)\cup (0;\infty )\)}{\((-\infty ;-3)\cup (3;\infty )\)}{\((-\infty ;-1)\cup (1;\infty )\)}{\((-1;1)\)}{\(\mathbb{R}\)}{}}
\LANGpl{
\Question{
Wyznacz zbiór rozwiązań podanej nierówności. 
\[ 
 (x + 2)(x^{2} + 4x + 3) > x^{2} + 5x + 6
\]}
{\((-3;-2)\cup (0;\infty )\)}{\((-\infty ;-3)\cup (3;\infty )\)}{\((-\infty ;-1)\cup (1;\infty )\)}{\((-1;1)\)}{\(\mathbb{R}\)}{}}
\LANGcs{
\Question{
Najděte řešení nerovnice \((x + 2)(x^{2} + 4x + 3) > x^{2} + 5x + 6\).}
{\((-3;-2)\cup (0;\infty )\)}{\((-\infty ;-3)\cup (3;\infty )\)}{\((-\infty ;-1)\cup (1;\infty )\)}{\((-1;1)\)}{\(\mathbb{R}\)}{}}
\LANGsk{
\Question{
Nájdite riešenie danej nerovnice. 
\[ 
 (x + 2)(x^{2} + 4x + 3) > x^{2} + 5x + 6
\]}
{\((-3;-2)\cup (0;\infty )\)}{\((-\infty ;-3)\cup (3;\infty )\)}{\((-\infty ;-1)\cup (1;\infty )\)}{\((-1;1)\)}{\(\mathbb{R}\)}{}}
\LANGes{
\Question{
Calcula el conjunto de soluciones de la siguiente inecuación.
\[ 
 (x + 2)(x^{2} + 4x + 3) > x^{2} + 5x + 6
\]}
{\((-3;-2)\cup (0;\infty )\)}{\((-\infty ;-3)\cup (3;\infty )\)}{\((-\infty ;-1)\cup (1;\infty )\)}{\((-1;1)\)}{\(\mathbb{R}\)}{}}
\End

\Data\ID{9000031001}
\Updated{2020-07-15 03:07:34}
\Level{B}
\SubArea{Higher degree equations and inequalities}
\LANGen{
\Question{
Find the sum of all real roots of the following equation. 
\[ 
 (3x - 1)(2x + 1)(4x^{2} + 3x - 1) = 0
\]}
{\(-\frac{11} 
{12}\)}{\(- \frac{1} 
{12}\)}{\(-\frac{1} 
{6}\)}{\(\frac{1}
{6}\)}{}{}}
\LANGpl{
\Question{
Oblicz sumę wszystkich rzeczywistych pierwiastków podanego równania.
\[ 
 (3x - 1)(2x + 1)(4x^{2} + 3x - 1) = 0
\]}
{\(-\frac{11} 
{12}\)}{\(- \frac{1} 
{12}\)}{\(-\frac{1} 
{6}\)}{\(\frac{1}
{6}\)}{}{}}
\LANGcs{
\Question{
Součet všech reálných kořenů rovnice 
\[ 
 (3x - 1)(2x + 1)(4x^{2} + 3x - 1) = 0
\]
je:}
{\(-\frac{11} 
{12}\)}{\(- \frac{1} 
{12}\)}{\(-\frac{1} 
{6}\)}{\(\frac{1}
{6}\)}{}{}}
\LANGsk{
\Question{
Nájdite súčet všetkých reálnych koreňov danej rovnice. 
\[ 
 (3x - 1)(2x + 1)(4x^{2} + 3x - 1) = 0
\]}
{\(-\frac{11} 
{12}\)}{\(- \frac{1} 
{12}\)}{\(-\frac{1} 
{6}\)}{\(\frac{1}
{6}\)}{}{}}
\LANGes{
\Question{
Calcula la suma de todas las raíces reales de la siguiente ecuación. 
\[ 
 (3x - 1)(2x + 1)(4x^{2} + 3x - 1) = 0
\]}
{\(-\frac{11} 
{12}\)}{\(- \frac{1} 
{12}\)}{\(-\frac{1} 
{6}\)}{\(\frac{1}
{6}\)}{}{}}
\End

\Data\ID{9000031002}
\Updated{2020-07-15 04:07:00}
\Level{B}
\SubArea{Higher degree equations and inequalities}
\LANGen{
\Question{
One of the solutions of the following equation is
\(x = 2\). Find
the set of all solutions. 
\[ 
 x^{3} + 2x^{2} - 5x - 6 = 0
\]}
{\(\{ - 3;-1;2\}\)}{\( \{ - 3;-1\}\)}{\( \{ - 3;0;2\}\)}{\(\{ - 1;2;3\}\)}{}{}}
\LANGpl{
\Question{
Jednym z rozwiązań podanego równania jest 
\(x = 2\). Wyznacz zbiór wszystkich rozwiązań.
\[ 
 x^{3} + 2x^{2} - 5x - 6 = 0
\]}
{\(\{ - 3;-1;2\}\)}{\(\{ - 3;-1\}\)}{\( \{ - 3;0;2\}\)}{\(\{ - 1;2;3\}\)}{}{}}
\LANGcs{
\Question{
Jestliže víme, že jeden z kořenů rovnice 
\[ 
 x^{3} + 2x^{2} - 5x - 6 = 0
\]
je roven \(2\),
pak množina všech kořenů této rovnice je:}
{\(\{ - 3;-1;2\}\)}{\( \{ - 3;-1\}\)}{\( \{ - 3;0;2\}\)}{\(\{ - 1;2;3\}\)}{}{}}
\LANGsk{
\Question{
Jeden z koreňov danej rovnice je 
\(x = 2\). Nájdite
množinu všetkých koreňov tejto rovnice. 
\[ 
 x^{3} + 2x^{2} - 5x - 6 = 0
\]}
{\(\{ - 3;-1;2\}\)}{\( \{ - 3;-1\}\)}{\( \{ - 3;0;2\}\)}{\(\{ - 1;2;3\}\)}{}{}}
\LANGes{
\Question{
Una de las soluciones de la siguiente ecuación es \(x = 2\). Calcula el conjunto de todas las soluciones. 
\[ 
 x^{3} + 2x^{2} - 5x - 6 = 0
\]}
{\(\{ - 3;-1;2\}\)}{\( \{ - 3;-1\}\)}{\( \{ - 3;0;2\}\)}{\(\{ - 1;2;3\}\)}{}{}}
\End

\Data\ID{9000031004}
\Updated{2020-07-15 03:07:34}
\Level{B}
\SubArea{Higher degree equations and inequalities}
\LANGen{
\Question{
Assuming \(y\in \mathbb{R}\),
find the number of the solutions of the following algebraic equation. 
\[ 
 y^{4} + 5y^{2} + 6 = 0
\]}
{\(0\)}{\(4\)}{\(3\)}{\(2\)}{}{}}
\LANGpl{
\Question{
Zakładając, że \(y\in \mathbb{R}\),
określ liczbę rozwiązań podanego wyrażenia algebraicznego.
\[ 
 y^{4} + 5y^{2} + 6 = 0
\]}
{\(0\)}{\(4\)}{\(3\)}{\(2\)}{}{}}
\LANGcs{
\Question{
Počet řešení rovnice 
\[ 
 y^{4} + 5y^{2} + 6 = 0
\]
v \(\mathbb{R}\) je:}
{\(0\)}{\(4\)}{\(3\)}{\(2\)}{}{}}
\LANGsk{
\Question{
Ak  \(y\in \mathbb{R}\),
nájdite počet riešení danej rovnice. 
\[ 
 y^{4} + 5y^{2} + 6 = 0
\]}
{\(0\)}{\(4\)}{\(3\)}{\(2\)}{}{}}
\LANGes{
\Question{
Suponiendo que \(y\in \mathbb{R}\),
encuentra el número de soluciones de la siguiente ecuación algebraica. 
\[ 
 y^{4} + 5y^{2} + 6 = 0
\]}
{\(0\)}{\(4\)}{\(3\)}{\(2\)}{}{}}
\End

\Data\ID{9000031005}
\Updated{2020-07-15 03:07:34}
\Level{B}
\SubArea{Higher degree equations and inequalities}
\LANGen{
\Question{
Assuming \(x\in \mathbb{R}\),
solve the following algebraic equation. 
\[ 
 (x + 1)^{4} - 5(x + 1)^{2} + 4 = 0
\]}
{\( \{ - 3;-2;0;1\}\)}{\( \{1;4\}\)}{\( \{ - 2;-1;1;2\}\)}{\( \{ - 1;3\}\)}{}{}}
\LANGpl{
\Question{
Zakładając, że \(x\in \mathbb{R}\),
rozwiąż podane równanie algebraiczne.
\[ 
 (x + 1)^{4} - 5(x + 1)^{2} + 4 = 0
\]}
{\( \{ - 3;-2;0;1\}\)}{\( \{1;4\}\)}{\( \{ - 2;-1;1;2\}\)}{\( \{ - 1;3\}\)}{}{}}
\LANGcs{
\Question{
Množinou všech řešení rovnice 
\[ 
 (x + 1)^{4} - 5(x + 1)^{2} + 4 = 0
\]
v \(\mathbb{R}\) je
množina:}
{\( \{ - 3;-2;0;1\}\)}{\( \{1;4\}\)}{\( \{ - 2;-1;1;2\}\)}{\( \{ - 1;3\}\)}{}{}}
\LANGsk{
\Question{
Ak \(x\in \mathbb{R}\),
vyriešte danú rovnicu. 
\[ 
 (x + 1)^{4} - 5(x + 1)^{2} + 4 = 0
\]}
{\( \{ - 3;-2;0;1\}\)}{\( \{1;4\}\)}{\( \{ - 2;-1;1;2\}\)}{\( \{ - 1;3\}\)}{}{}}
\LANGes{
\Question{
Suponiendo que \(x\in \mathbb{R}\), resuelve la siguiente ecuación algebraica. 
\[ 
 (x + 1)^{4} - 5(x + 1)^{2} + 4 = 0
\]}
{\( \{ - 3;-2;0;1\}\)}{\( \{1;4\}\)}{\( \{ - 2;-1;1;2\}\)}{\( \{ - 1;3\}\)}{}{}}
\End

\Data\ID{9000031008}
\Updated{2020-07-15 03:07:34}
\Level{B}
\SubArea{Higher degree equations and inequalities}
\LANGen{
\Question{
Assuming \(x\in \mathbb{R}\),
solve the following equation. 
\[ 
 4x^{3} - 3x^{2} - x = 0
\]}
{\( \left \{-\frac{1}
{4};0;1\right \}\)}{\(\{0;1;4\}\)}{\( \{1;4\}\)}{\( \{0\}\)}{}{}}
\LANGpl{
\Question{
Zakładając, że \(x\in \mathbb{R}\),
rozwiąż podane równanie.
\[ 
 4x^{3} - 3x^{2} - x = 0
\]}
{\( \left \{-\frac{1}
{4};0;1\right \}\)}{\( \{0;1;4\}\)}{\( \{1;4\}\)}{\( \{0\}\)}{}{}}
\LANGcs{
\Question{
Množinou všech řešení rovnice 
\[ 
 4x^{3} - 3x^{2} - x = 0
\]
v \(\mathbb{R}\) je
množina:}
{\( \left \{-\frac{1}
{4};0;1\right \}\)}{\( \{0;1;4\}\)}{\( \{1;4\}\)}{\( \{0\}\)}{}{}}
\LANGsk{
\Question{
Ak \(x\in \mathbb{R}\),
vyriešte danú rovnicu. 
\[ 
 4x^{3} - 3x^{2} - x = 0
\]}
{\(\left \{-\frac{1}
{4};0;1\right \}\)}{\( \{0;1;4\}\)}{\( \{1;4\}\)}{\( \{0\}\)}{}{}}
\LANGes{
\Question{
Suponiendo que \(x\in \mathbb{R}\), resuelve la ecuación siguiente. 
\[ 
 4x^{3} - 3x^{2} - x = 0
\]}
{\( \left \{-\frac{1}
{4};0;1\right \}\)}{\(\{0;1;4\}\)}{\( \{1;4\}\)}{\( \{0\}\)}{}{}}
\End

\Data\ID{9000031009}
\Updated{2020-07-15 03:07:34}
\Level{B}
\SubArea{Higher degree equations and inequalities}
\LANGen{
\Question{
Find the sum of the solutions of the following equation. 
\[ 
 6(3x + 1)(2x^{2} + 3x - 2) = 0
\]}
{\(-\frac{11} 
 {6} \)}{\(-\frac{7} 
{6}\)}{\(-\frac{1} 
{2}\)}{\(\frac{11}
 {6} \)}{}{}}
\LANGpl{
\Question{
Wyznacz sumę rozwiązań podanego równania.
\[ 
 6(3x + 1)(2x^{2} + 3x - 2) = 0
\]}
{\(-\frac{11} 
 {6} \)}{\(-\frac{7} 
{6}\)}{\(-\frac{1} 
{2}\)}{\(\frac{11}
 {6} \)}{}{}}
\LANGcs{
\Question{
Součet kořenů rovnice 
\[ 
 6(3x + 1)(2x^{2} + 3x - 2) = 0
\]
je roven:}
{\(-\frac{11} 
 {6} \)}{\(-\frac{7} 
{6}\)}{\(-\frac{1} 
{2}\)}{\(\frac{11}
 {6} \)}{}{}}
\LANGsk{
\Question{
Určte súčet koreňov danej rovnice. 
\[ 
 6(3x + 1)(2x^{2} + 3x - 2) = 0
\]}
{\(-\frac{11} 
 {6} \)}{\(-\frac{7} 
{6}\)}{\(-\frac{1} 
{2}\)}{\(\frac{11}
 {6} \)}{}{}}
\LANGes{
\Question{
Calcula la suma de las soluciones de la siguiente ecuación. 
\[ 
 6(3x + 1)(2x^{2} + 3x - 2) = 0
\]}
{\(-\frac{11} 
 {6} \)}{\(-\frac{7} 
{6}\)}{\(-\frac{1} 
{2}\)}{\(\frac{11}
 {6} \)}{}{}}
\End

\Data\ID{9000031010}
\Updated{2020-07-15 03:07:34}
\Level{B}
\SubArea{Higher degree equations and inequalities}
\LANGen{
\Question{
Identify a true statement on the following equation. 
\[ 
 x^{5} - x^{3} - 6x = 0
\]}
{The equation has three solutions in \(\mathbb{R}\).}{The equation does not have solution in
\(\mathbb{R}\).}{The equation has five solutions in \(\mathbb{R}\).}{The equation has one solution in \(\mathbb{R}\).}{}{}}
\LANGpl{
\Question{
Które z poniższych zdań dotyczących podanego równania jest prawdziwe?
\[ 
 x^{5} - x^{3} - 6x = 0
\]}
{Równanie ma trzy rozwiązania w  \(\mathbb{R}\).}{Równanie nie ma rozwiązania w 
\(\mathbb{R}\).}{Równanie ma pięć rozwiązań w \(\mathbb{R}\).}{Równanie ma jedno rozwiązanie w  \(\mathbb{R}\).}{}{}}
\LANGcs{
\Question{
Rovnice 
\[ 
 x^{5} - x^{3} - 6x = 0
\]}
{má právě \(3\) 
řešení v \(\mathbb{R}\).}{nemá v \(\mathbb{R}\) 
žádné řešení.}{má právě \(5\) 
řešení v \(\mathbb{R}\).}{má právě jedno řešení v \(\mathbb{R}\).}{}{}}
\LANGsk{
\Question{
Rovnica 
\[ 
 x^{5} - x^{3} - 6x = 0
\]}
{má práve tri riešenia v \(\mathbb{R}\).}{nemá žiadne riešenie v
\(\mathbb{R}\).}{má práve päť riešení v \(\mathbb{R}\).}{má práve jedno riešenie v \(\mathbb{R}\).}{}{}}
\LANGes{
\Question{
Identifica la proposición verdadera respecto a la siguiente ecuación. 
\[ 
 x^{5} - x^{3} - 6x = 0
\]}
{La ecuación tiene tres soluciones en \(\mathbb{R}\).}{La ecuación no tiene solución en
\(\mathbb{R}\).}{La ecuación tiene cinco soluciones en \(\mathbb{R}\).}{La ecuación tiene una solución en \(\mathbb{R}\).}{}{}}
\End

\Data\ID{1003029005}
\Updated{2020-07-15 03:07:12}
\Level{C}
\SubArea{Higher degree equations and inequalities}
\LANGen{
\Question{
Find the solution set  of the inequality.
\[ \left(2x^2-18\right)\left(x^3-8\right) < 0 \]}
{\( (-\infty;-3)\cup(2;3) \)}{\( (-3;2)\cup(3;\infty) \)}{\( (-3;3) \)}{\( (-\infty;-3) \)}{}{}}
\LANGpl{
\Question{
Wyznacz zbiór rozwiązań podanej nierówności.
\[ \left(2x^2-18\right)\left(x^3-8\right) < 0 \]}
{\( (-\infty;-3)\cup(2;3) \)}{\( (-3;2)\cup(3;\infty) \)}{\( (-3;3) \)}{\( (-\infty;-3) \)}{}{}}
\LANGcs{
\Question{
Určete množinu řešení nerovnice.
\[ \left(2x^2-18\right)\left(x^3-8\right) < 0 \]}
{\( (-\infty;-3)\cup(2;3) \)}{\( (-3;2)\cup(3;\infty) \)}{\( (-3;3) \)}{\( (-\infty;-3) \)}{}{}}
\LANGsk{
\Question{
Určte množinu riešení nerovnice.
\[ \left(2x^2-18\right)\left(x^3-8\right) < 0 \]}
{\( (-\infty;-3)\cup(2;3) \)}{\( (-3;2)\cup(3;\infty) \)}{\( (-3;3) \)}{\( (-\infty;-3) \)}{}{}}
\LANGes{
\Question{
Calcula el conjunto de soluciones de la inecuación.
\[ \left(2x^2-18\right)\left(x^3-8\right) < 0 \]}
{\( (-\infty;-3)\cup(2;3) \)}{\( (-3;2)\cup(3;\infty) \)}{\( (-3;3) \)}{\( (-\infty;-3) \)}{}{}}
\End

\Data\ID{9000028304}
\Updated{2020-07-15 03:07:34}
\Level{C}
\SubArea{Higher degree equations and inequalities}
\LANGen{
\Question{
The following equation has solutions
\(x_1= 1\) and
\(x_2 = 3\). Find
the sum of the remaining real solutions. 
\[ 
 x^{4} - 12x^{3} + 47x^{2} - 72x + 36 = 0
\]}
{\(8\)}{\(- 1\)}{\(3\)}{\(5\)}{}{}}
\LANGpl{
\Question{
Rozwiązaniem podanego równania jest 
\(x_1 = 1\) i
\(x_2 = 3\). Oblicz sumę pozostałych rzeczywistych rozwiązań.
\[ 
 x^{4} - 12x^{3} + 47x^{2} - 72x + 36 = 0
\]}
{\(8\)}{\(- 1\)}{\(3\)}{\(5\)}{}{}}
\LANGcs{
\Question{
Daná rovnice má dva kořeny $x_1=1$ a $x_2=2$. Určete součet zbývajících reálných kořenů rovnice.
\[x^{4} - 12x^{3} + 47x^{2} - 72x + 36 = 0\]}
{\(8\)}{\(- 1\)}{\(3\)}{\(5\)}{}{}}
\LANGsk{
\Question{
Daná rovnica má dva korene 
\(x_1 = 1\) a
\(x_2 = 3\). Určte súčet zvyšných reálnych koreňov rovnice.
\[ 
 x^{4} - 12x^{3} + 47x^{2} - 72x + 36 = 0
\]}
{\(8\)}{\(- 1\)}{\(3\)}{\(5\)}{}{}}
\LANGes{
\Question{
La siguiente ecuación tiene  soluciones \(x_1= 1\) y \(x_2 = 3\). Calcula la suma de las soluciones reales restantes.
\[ 
 x^{4} - 12x^{3} + 47x^{2} - 72x + 36 = 0
\]}
{\(8\)}{\(- 1\)}{\(3\)}{\(5\)}{}{}}
\End

\Data\ID{9000028305}
\Updated{2020-07-15 03:07:34}
\Level{C}
\SubArea{Higher degree equations and inequalities}
\LANGen{
\Question{
The following equation has solution
\(x_1= 2\) and
\(x_2= 4\). Find
the sum of the remaining real solutions. 
\[ 
 x^{4} - 6x^{3} - x^{2} + 54x - 72 = 0
\]}
{\(0\)}{\(- 1\)}{\(1\)}{\(2\)}{}{}}
\LANGpl{
\Question{
Rozwiązaniem podanego równania jest 
\(x_1 = 2\) i 
\(x_2 = 4\). Oblicz sumę pozostałych rzeczywistych rozwiązań.
\[ 
 x^{4} - 6x^{3} - x^{2} + 54x - 72 = 0
\]}
{\(0\)}{\(- 1\)}{\(1\)}{\(2\)}{}{}}
\LANGcs{
\Question{
Daná rovnice má dva kořeny $x_1=1$ a $x_2=4$. Určete součet zbývajících reálných kořenů rovnice.
\[x^{4} - 6x^{3} - x^{2} + 54x - 72 = 0\]}
{\(0\)}{\(- 1\)}{\(1\)}{\(2\)}{}{}}
\LANGsk{
\Question{
Daná rovnica má dva korene 
\(x_1 = 2\) a
\(x_2 = 4\). Určte súčet zvyšných  reálnych koreňov rovnice.
\[ 
 x^{4} - 6x^{3} - x^{2} + 54x - 72 = 0
\]}
{\(0\)}{\(- 1\)}{\(1\)}{\(2\)}{}{}}
\LANGes{
\Question{
La siguiente ecuación tiene soluciones \(x_1= 2\) y \(x_2= 4\). Calcula la suma de las soluciones reales restantes. 
\[ 
 x^{4} - 6x^{3} - x^{2} + 54x - 72 = 0
\]}
{\(0\)}{\(- 1\)}{\(1\)}{\(2\)}{}{}}
\End

\Data\ID{9000028309}
\Updated{2020-07-15 03:07:34}
\Level{C}
\SubArea{Higher degree equations and inequalities}
\LANGen{
\Question{
Find the sum of all real solutions of the following equation. 
\[ 
 x^{4} + x^{3} + x^{2} + x = 0
\]}
{\(- 1\)}{\(0\)}{\(5\)}{\(6\)}{}{}}
\LANGpl{
\Question{
Oblicz sumę wszystkich rzeczywistych rozwiązań podanego równania.
\[ 
 x^{4} + x^{3} + x^{2} + x = 0
\]}
{\(- 1\)}{\(0\)}{\(5\)}{\(6\)}{}{}}
\LANGcs{
\Question{
Určete součet všech reálných kořenů dané rovnice.
\[ 
 x^{4} + x^{3} + x^{2} + x = 0
\]}
{\(- 1\)}{\(0\)}{\(5\)}{\(6\)}{}{}}
\LANGsk{
\Question{
Určte súčet všetkých reálnych koreňov danej rovnice. 
\[ 
 x^{4} + x^{3} + x^{2} + x = 0
\]}
{\(- 1\)}{\(0\)}{\(5\)}{\(6\)}{}{}}
\LANGes{
\Question{
Calcula la suma de todas las soluciones reales de la siguiente ecuación. 
\[ 
 x^{4} + x^{3} + x^{2} + x = 0
\]}
{\(- 1\)}{\(0\)}{\(5\)}{\(6\)}{}{}}
\End

\Data\ID{9000029306}
\Updated{2020-07-15 03:07:34}
\Level{C}
\SubArea{Higher degree equations and inequalities}
\LANGen{
\Question{
Find the solution set of the following inequality. 
\[ 
 x^{3} - 3x^{2} + 3x - 1 < 0
\]}
{\(\left (-\infty ;1\right )\)}{\(\emptyset \)}{\(\mathbb{R}\)}{\(\left (3;\infty \right )\)}{}{}}
\LANGpl{
\Question{
Wyznacz zbiór rozwiązań podanej nierówności. 
\[ 
 x^{3} - 3x^{2} + 3x - 1 < 0
\]}
{\(\left (-\infty ;1\right )\)}{\(\emptyset \)}{\(\mathbb{R}\)}{\(\left (3;\infty \right )\)}{}{}}
\LANGcs{
\Question{
Vyberte množinu řešení následující nerovnice. \[x^{3} - 3x^{2} + 3x - 1 < 0\]}
{\(\left (-\infty ;1\right )\)}{$\emptyset$}{\(\mathbb{R}\)}{\(\left (3;\infty \right )\)}{}{}}
\LANGsk{
\Question{
Nájdite riešenie danej nerovnice. 
\[ 
 x^{3} - 3x^{2} + 3x - 1 < 0
\]}
{\(\left (-\infty ;1\right )\)}{\(\emptyset \)}{\(\mathbb{R}\)}{\(\left (3;\infty \right )\)}{}{}}
\LANGes{
\Question{
Calcula el conjunto de soluciones de la siguiente inecuación.. 
\[ 
 x^{3} - 3x^{2} + 3x - 1 < 0
\]}
{\(\left (-\infty ;1\right )\)}{\(\emptyset \)}{\(\mathbb{R}\)}{\(\left (3;\infty \right )\)}{}{}}
\End

\Data\ID{9000031003}
\Updated{2020-07-15 03:07:34}
\Level{C}
\SubArea{Higher degree equations and inequalities}
\LANGen{
\Question{
Assuming \(x\in \mathbb{R}\),
solve the following algebraic equation. 
\[ 
 x^{4} + 4x^{2} - 5 = 0
\]}
{\( \{ - 1;1\}\)}{\( \{1\}\)}{\( \{ -\sqrt{5};-1;1;\sqrt{5}\}\)}{\( \emptyset \)}{}{}}
\LANGpl{
\Question{
Zakładając, że \(x\in \mathbb{R}\), rozwiąż podane równanie algebraiczne.
\[ 
 x^{4} + 4x^{2} - 5 = 0
\]}
{\( \{ - 1;1\}\)}{\( \{1\}\)}{\( \{ -\sqrt{5};-1;1;\sqrt{5}\}\)}{\( \emptyset \)}{}{}}
\LANGcs{
\Question{
Množinou všech řešení rovnice 
\[ 
 x^{4} + 4x^{2} - 5 = 0
\]
v \(\mathbb{R}\) je
množina:}
{\(\{ - 1;1\}\)}{\( \{1\}\)}{\( \{ -\sqrt{5};-1;1;\sqrt{5}\}\)}{\( \emptyset \)}{}{}}
\LANGsk{
\Question{
Ak \(x\in \mathbb{R}\),
vyriešte danú rovnicu. 
\[ 
 x^{4} + 4x^{2} - 5 = 0
\]}
{\( \{ - 1;1\}\)}{\( \{1\}\)}{\( \{ -\sqrt{5};-1;1;\sqrt{5}\}\)}{\( \emptyset \)}{}{}}
\LANGes{
\Question{
Suponiendo que \(x\in \mathbb{R}\), resuelve la siguiente ecuación algebraica. 
\[ 
 x^{4} + 4x^{2} - 5 = 0
\]}
{\( \{ - 1;1\}\)}{\( \{1\}\)}{\( \{ -\sqrt{5};-1;1;\sqrt{5}\}\)}{\( \emptyset \)}{}{}}
\End

\Data\ID{9000031006}
\Updated{2020-07-15 03:07:34}
\Level{C}
\SubArea{Higher degree equations and inequalities}
\LANGen{
\Question{
The following equation has a double solution
\(x = 1\). Find
all solutions. 
\[ 
 x^{4} + 2x^{3} - 3x^{2} - 4x + 4 = 0
\]}
{\(K = \{ - 2;1\}\)}{\(K = \{ - 2;1;2\}\)}{\(K = \{ - 2;0;1\}\)}{another answer}{}{}}
\LANGpl{
\Question{
Podane równanie ma podwójne rozwiązanie
\(x = 1\). Znajdź wszystkie rozwiązania.
\[ 
 x^{4} + 2x^{3} - 3x^{2} - 4x + 4 = 0
\]}
{\(K = \{ - 2;1\}\)}{\(K = \{ - 2;1;2\}\)}{\(K = \{ - 2;0;1\}\)}{inna odpowiedź}{}{}}
\LANGcs{
\Question{
Víte-li, že jeden dvojnásobný kořen rovnice 
\[ 
 x^{4} + 2x^{3} - 3x^{2} - 4x + 4 = 0
\]
je roven \(1\),
pak množinou všech kořenů této rovnice je:}
{\(K = \{ - 2;1\}\)}{\(K = \{ - 2;1;2\}\)}{\(K = \{ - 2;0;1\}\)}{žádná z uvedených množin}{}{}}
\LANGsk{
\Question{
Daná rovnica má dvojnásobný koreň 
\(x = 1\). Určte množinu všetkých koreňov tejto rovnice.
\[ 
 x^{4} + 2x^{3} - 3x^{2} - 4x + 4 = 0
\]}
{\(K = \{ - 2;1\}\)}{\(K = \{ - 2;1;2\}\)}{\(K = \{ - 2;0;1\}\)}{žiadna z uvedených množín}{}{}}
\LANGes{
\Question{
La siguiente ecuación tiene una solución doble  \(x = 1\). Calcula todas las soluciones.
\[ 
 x^{4} + 2x^{3} - 3x^{2} - 4x + 4 = 0
\]}
{\(K = \{ - 2;1\}\)}{\(K = \{ - 2;1;2\}\)}{\(K = \{ - 2;0;1\}\)}{another answer}{}{}}
\End

\Data\ID{9000031007}
\Updated{2020-07-15 03:07:34}
\Level{C}
\SubArea{Higher degree equations and inequalities}
\LANGen{
\Question{
Find the sum of the solutions of the following equation. 
\[ 
 x^{3} + 2x^{2} - x - 2 = 0
\]}
{\(- 2\)}{\(3\)}{\(- 3\)}{\(- 1\)}{}{}}
\LANGpl{
\Question{
Oblicz sumę rozwiązań podanego równania.
\[ 
 x^{3} + 2x^{2} - x - 2 = 0
\]}
{\(- 2\)}{\(3\)}{\(- 3\)}{\(- 1\)}{}{}}
\LANGcs{
\Question{
Součet kořenů rovnice 
\[ 
 x^{3} + 2x^{2} - x - 2 = 0
\]
je roven:}
{\(- 2\)}{\(3\)}{\(- 3\)}{\(- 1\)}{}{}}
\LANGsk{
\Question{
Určte súčet koreňov danej rovnice. 
\[ 
 x^{3} + 2x^{2} - x - 2 = 0
\]}
{\(- 2\)}{\(3\)}{\(- 3\)}{\(- 1\)}{}{}}
\LANGes{
\Question{
Calcula la suma de las soluciones de la siguiente ecuación. 
\[ 
 x^{3} + 2x^{2} - x - 2 = 0
\]}
{\(- 2\)}{\(3\)}{\(- 3\)}{\(- 1\)}{}{}}
\End
